\def\topic{}
\ifdefined\MyHeaderIncluded
\else
\def\qwertyflag{true}

\ifdefined\MyHeaderIncluded

\else
  \newcommand{\foobar}{FooBar}
  \def\MyHeaderIncluded{true}
% if is till the end of the file

\documentclass[10pt,a4paper]{article}
\usepackage[utf8]{inputenc}
%\usepackage[german]{babel}
%\usepackage[cyrillic]{babel}
\usepackage[T1]{fontenc}
%\usepackage[T2A]{fontenc}
%\usepackage[english, russian]{babel}
\usepackage{amsmath}
\usepackage{amsfonts}
\usepackage{amssymb}
\usepackage{amsthm}
\usepackage{mathtools} %coloneqq

\usepackage{bussproofs} %

%\renewcommand{\thesubsection}{\arabic{subsection}}
%\newtheorem{thm}{Theorem}[section]
%\renewcommand{\thethm}{\arabic{subsection}.\arabic{thm}}

\usepackage{thmtools}
\usepackage{hyperref}
\hypersetup{
    colorlinks=true,
    linkcolor=blue,
    filecolor=magenta,      
    urlcolor=cyan,
    pdftitle={Overleaf Example},
    %pdfpagemode=FullScreen,
    }
%\usepackage{stix-type1}%\turnediota
\theoremstyle{definition}

\newtheorem{constructor}{Constructor}[section]
\newtheorem{iaxiom}{Axiom}
\newtheorem{axiom}{Axiom}
\newtheorem*{axiom*}{Axiom}
\newtheorem{definition}{Definition}[section]
\newtheorem*{definition*}{Definition}
\newtheorem{remark}{Remark}[section]
\newtheorem*{remark*}{Remark}
\newtheorem{recall}{Recall}[section]
\newtheorem*{recall*}{Recall}

\newtheorem{lemma}{Lemma}[section]
\newtheorem{theorem}{Theorem}[section]
\newtheorem*{theorem*}{Theorem}
\newtheorem{note}{Note}[section]
\usepackage{etoolbox}

%\newcommand*{\@Dashv}[2]{%
%  \reflectbox{$\m@th#1#2$}%
%}
\newcommand*{\separation}[3]{% 1)binder, 2)set, 3)condition
  \{#1\in #2: #3\}
}

% https://www.freecodecamp.org/news/simple-regex-tricks-for-beginners-3acb3fa257cb/
%  \\mathcal\{T\} |---> \topology
%  \\mathcal\{([^\}]*)\}   |--->  $1
\newcommand{\AND}{\ \ \&\ \ }
\newcommand{\any}{\mathrm{\ any}}
\newcommand{\borel}{\mathcal{B}}
\newcommand{\can}{\mathrm{can}}
\newcommand{\canonize}{\mathrm{canonize}}
\newcommand{\closure}{\mbox{cl}}
\newcommand{\complex}{\mathbb{C}}
\newcommand{\converges}{\longrightarrow}
\newcommand{\countable}{\mbox{\ countable}}
\newcommand{\ldefined}{\mbox{\ defined}}
\newcommand{\dom}{\mbox{dom}}
\newcommand{\Dom}{\mbox{Dom}}
\newcommand{\Done}{\noindent \textbf{Done.}\\}
\newcommand{\lexists}{\mbox{\ exists}}
\newcommand{\lforall}{\mbox{\ forall}}
\newcommand{\lunique}{\mbox{\ unique}}
\newcommand{\family}{\mathrm{\ family}}
\newcommand{\Fm}{\mathrm{Fm}}
\newcommand{\FuncSymb}{\mathrm{FuncSymb}}
\newcommand{\function}{\mathrm{\ function}}
\newcommand{\functional}{\mbox{\ functional}}
\newcommand{\IFF}{\ \Longleftrightarrow\ }
\renewcommand{\iff}{\ \longleftrightarrow\ }
\newcommand{\IMPLIES}{\ \Longrightarrow\ }
\renewcommand{\implies}{\ \longrightarrow\ }
\newcommand{\Inf}{\mathrm{Inf}}
\newcommand{\injective}{\mbox{\ injective}}
\newcommand{\injection}{\mbox{\ injection}}
\newcommand{\interior}{\mbox{int}}
\newcommand{\intersection}{\bigcap}
\newcommand{\isempty}{\mathrm{\ empty}}
\newcommand{\LB}{\mathrm{LB}}
\newcommand{\nghbset}{\mathcal{N}}
\newcommand{\nonempty}{\mathrm{\ nonempty}}
\newcommand{\Ord}{\mathrm{Ord}}
\newcommand{\power}{\mathcal{P}}
\newcommand{\pr}{\mathrm{pr}}
\newcommand{\fst}{\mathrm{fst}}
\newcommand{\snd}{\mathrm{snd}}
\newcommand{\PredSymb}{\mathrm{PredSymb}}
\newcommand{\Proof}{\noindent \textbf{Proof.}\\}
\newcommand{\ran}{\mathrm{ran}}
\newcommand{\Ran}{\mathrm{Ran}}
\newcommand{\real}{\mathbb{R}}
\newcommand{\relation}{\mathrm{\ relation}}
\newcommand{\set}{\mathrm{\ set}}
\newcommand{\singleton}{\mathrm{\ singleton}}
\newcommand{\semantics}{\mathrm{semantics}}
\newcommand{\semialgebra}{\mathcal{S}}
\newcommand{\Sup}{\mathrm{Sup}}
\newcommand{\surjective}{\mbox{\ surjective\ on\ }}
\newcommand{\Term}{\mathrm{Term}}
\newcommand{\Type}{\mathrm{Type}}
\newcommand{\topology}{\mathcal{T}}
\newcommand{\total}{\mbox{\ total\ on\ }}
\newcommand{\transitive}{\mathrm{\ transitive}}
\newcommand{\union}{\bigcup}
\newcommand{\lTrue}{\mathrm{True}}
\newcommand{\lFalse}{\mathrm{False}}
\newcommand{\lUncertain}{\mathrm{Uncertain}}
\newcommand{\lif}{\mathrm{if}}
\newcommand{\lthen}{\mathrm{then}}
\newcommand{\lelse}{\mathrm{else}}
\newcommand{\lexception}{\mathrm{exception}}
\newcommand{\lundef}{\mathrm{undef}}
\newcommand{\univocal}{\mathrm{\ univocal}}
\newcommand{\univocalat}{\mathrm{\ univocal\ at\ }}
\newcommand{\UB}{\mathrm{UB}}
\newcommand{\Un}{\mathcal{U}}
\newcommand{\Var}{\mathrm{Var}}

\newcommand{\notvarphi}{\phi}
\renewcommand{\phi}{\varphi}

\newcommand{\comment}[1]{\quad\mbox{(#1)}}
% local commands
%\preto\section{\undef\blipo} %etoolbox
\preto\subsection{\undef\blipo} %etoolbox
% https://tex.stackexchange.com/questions/105684/limit-macro-newcommand-scope-to-section
\usepackage{graphicx}
\usepackage{authblk}

\setlength\parindent{0pt} % whole doc of \noindent
\usepackage{soul}
%\usepackage{mathgl}

\fi % guard
\ifdefined\topic
\else
\def\topic{No Titel}
\fi
\title{\topic}
\author[1]{Georgy Dunaev}
\begin{document}
\maketitle
\fi
\section{Types as predicates on classes}

%\section{[OUTDATED] Classes--as--types, sets--as--expressions approach}

%\subsection{Classes--as--Types semantics}
This section investigates interpretations of dependent sums, dependent
products, and inductive and coinductive types using classes. Its central
idea is that a type specifies a unary predicate on possible term values.
The instances may include proper classes, so their collection need not
itself be a class in the object theory.

\paragraph{Local models and earlier proposals.}
The subsections compare several representations; they do not assert that
all displayed models are interchangeable. In particular, a class of set
values, a predicate on arbitrary classes, and a syntactic type expression
belong to different levels. A soundness claim must identify which of these
is being interpreted and which family encoding it uses.

\begin{note}[OUTDATED: the unrestricted class-as-type proposal]
An earlier proposal assigned truth to every type declaration:
\[
\semantics(``(A:\Type)")\eqdef\top.
\]
This does not enforce the intended uniqueness of the term--type relation.
It is retained as motivation for interpreting types by explicit predicates.
In that earlier family-based proposal, a dependent declaration was written
\[
\semantics(``(x:A\vdash B(x):\Type)")
\eqdef (B\family)\land(\Dom(B)\subseteq A),
\]
and the semisingleton encoding was
\[
(\Lambda x,\ t)=\cc{(x,s)}{(s\subseteq t)\land(s\ \mbox{ss})}.
\]
The domain condition and evaluation rules for a dependent type must be
specified in the particular model; this sketch does not supply them all.
\end{note}

\subsection{Explicit typed term interpretation ($\mu_A$ / $I$)}
Throughout this chapter, the background model of the class theory (NBG/MK) is fixed (though unspecified), and the valuation of free variables is likewise fixed (though unspecified).
With that context fixed, the interpretation map $\mu_A$ (written $I(a)$ or $\mu_A(a)$ below, and rendered as \texttt{\textbackslash mu\_A} in the Obelisk SML kernel \texttt{typetheory.ML}) is an \emph{explicit typed term interpretation} (an expression-level interpretation map):
\[
\mu_A : A.\mathsf{Expression} \to \mathsf{Class}.\mathsf{Expression},
\]
which translates a typed expression $a:A$ into a class expression $\mu_A(a)\equiv I(a)$. Evaluating the translated expression in the fixed model is a further semantic operation; the expression itself is not a class object. The term \emph{valuation homomorphism}
can be used for the resulting semantic interpretation when preservation of
the relevant constructors has been proved. That preservation is an
obligation of the model, not a consequence of the map's displayed type.

\begin{note}[advanced semantic setting: terms as classes, conglomerates (= predicates on classes) as types]
Unlike the standard set-theoretic setting where terms are interpreted as sets and classes are interpreted as types, we work in a one-level-higher setting:
\begin{enumerate}
\item \textbf{Terms are interpreted as classes} (which may be either sets or proper classes such as $\V$ and $\On$).
\item \textbf{Types are interpreted as conglomerates}, which we use interchangeably with \textbf{predicates on the conglomerate of all classes} ($\Class \to \Bool$). In particular, the type $\Class$ is the constant-true predicate on classes $\lambda (C:\Class).\ \top$ (informally, a constant conglomerate-function to $\Bool$, representing the conglomerate of all classes).
\item \textbf{$\Class$ is the default type}: it is implicitly assumed whenever a quantifier or comprehension omits an explicit type ($\forall x,\ \varphi(x) :\equiv \forall (x:\Class),\ \varphi(x)$) and in \textbf{bound (bounded) quantifiers}:
\[
\begin{aligned}
\forall (x\in Y),\ \varphi(x) &\iffdef \forall (x:\Class),\ (x\in Y)\implies \varphi(x),\\
\exists (x\in Y),\ \varphi(x) &\iffdef \exists (x:\Class),\ (x\in Y)\land \varphi(x).
\end{aligned}
\]
\end{enumerate}
\end{note}

\subsection{Storage [unfinished]}
Computer memory stores finite representations. A function expression can
therefore be stored as a bound variable together with a term, rather than
as a list of all its argument--value pairs. Types rule out expressions
that do not fit the intended interfaces. A semantic interpretation then
provides proof obligations for the implementation: the rules must preserve
the interpretation, and a failed obligation may reveal an implementation
error or a missing hypothesis.

\iffalse
Is the following okey ???
\[
(B\family)\iff\exists t : B = (\Lambda x: t)
\]
\fi
\pagebreak
\subsection{Implementing the boolean type}

\subsubsection{Axioms of the boolean type}
The class-valued conditional below has two computation laws. We then
give a set model of the Boolean type with its dependent elimination rule.
These are related constructions: a Boolean value is an element of a
two-element set, whereas a proposition is a formula.

\subsubsection{Model of boolean type}
\underline{Formation of boolean type:}\\
\underline{Introduction of terms of boolean type:}\\
\underline{Elimination of terms of boolean type:}\\
\begin{definition}[conditional operator if--then--else / $\lifdep$]
\[
(\lif\ \varphi\ \lthen\ t_1\ \lelse\ t_2) \eqdef \lifdep(\varphi, t_1, t_2) \eqdef (\cc{z}{((z\in t_1)\land\varphi)\lor((z\in t_2)\land\neg\varphi)})
\]
Side conditions:\\
($z$ is not free in $t_1$),\\
($z$ is not free in $t_2$),\\
($z$ is not free in $\varphi$).\\
\end{definition}
\underline{Computation of terms of boolean type:}\\
\begin{theorem}[$\lifdep$--true--elim]
\[
(\varphi \implies (\lifdep(\varphi, t_1, t_2) = t_1))
\]
\end{theorem}
\begin{theorem}[$\lifdep$--false--elim]
\[
((\neg\varphi) \implies (\lifdep(\varphi, t_1, t_2) = t_2))
\]
\end{theorem}

\subsubsection{Soundness of the boolean type}
Under $\varphi$, the comprehension defining $\lifdep$ has exactly the
members of $t_1$; under $\neg\varphi$, it has exactly the members of $t_2$.
Extensionality therefore proves the displayed computation laws. Each
conditional implication holds intuitionistically; selecting a branch for
an arbitrary formula requires deciding that formula.



\begin{definition}[The two-element set model]
Put $0=\emptyclass$, $1=\{0\}$ and
$\mathrm{Bool}=\{0,1\}$. Interpret $\mathsf{false}$ by $0$ and
$\mathsf{true}$ by $1$. Pairing and the empty-set axiom make these sets.
They are distinct because $0\in1$ but $0$ has no members. Thus every
$b\in\mathrm{Bool}$ is exactly one of these constructors.
\end{definition}

\begin{theorem}[Dependent Boolean elimination]
Let $E$ be a set-valued family on $\mathrm{Bool}$ and suppose
$e_0\in E(0)$, $e_1\in E(1)$. There is a unique function $r$ with
domain $\mathrm{Bool}$ such that $r(b)\in E(b)$ and
\[
r(0)=e_0,\qquad r(1)=e_1.
\]
For every dependent function $q$ on $\mathrm{Bool}$, elimination with
branches $q(0)$ and $q(1)$ returns $q$ itself.
\end{theorem}
\Proof
Use the graph $\{(0,e_0),(1,e_1)\}$ with ordinary injective set pairs.
Distinctness of $0,1$ makes it functional with precisely the desired
domain, and its values have the required types by hypothesis.
Any other function satisfying the equations agrees at both possible
arguments and hence equals it by extensionality. Taking $e_i=q(i)$
proves the last assertion, the extensional Boolean eta law.
\Done

\begin{corollary}[Boolean induction]
If a predicate $P$ holds at $0$ and $1$, it holds at every element of
$\mathrm{Bool}$.
\end{corollary}
\Proof
Membership in $\{0,1\}$ gives $b=0$ or $b=1$. Substitute in the
corresponding premise. This uses the definition of the pair set, not a
law deciding every proposition in the surrounding logic.
\Done

For a decidable formula $\phi$, the Boolean characteristic value is $1$
when $\phi$ and $0$ when $\neg\phi$. Applying elimination with branches
$t_2,t_1$ agrees with the earlier conditional when these branch values
are sets. That agreement follows by the two cases. Arbitrary class-valued
branches still have the earlier comprehension-based conditional, but they
cannot both be stored as values of an ordinary set graph unless they are
sets. This is the same size distinction encountered in family encodings.

\subsection{Implementing option types}
The concrete semisingleton model below concerns set-valued instances of
$A$: the empty class represents no value, and a singleton represents one
value. It does not directly encode a proper-class value as a singleton.
The draft uses both $\Emb(A)$ as a predicate and its collected extension
$\kappa(\Emb(A))$; these must be distinguished whenever membership notation
is used. The displayed typing rules record the intended interface while
the set-class calculations establish the concrete model's basic cases.

%Axioms:\\
\subsubsection{Axioms of option types}
\begin{axiom}[None--ax]
\[
\None:\Option(A)
\]
\[
I(\None)\in\kappa(\Emb(\Option(A)))
\]
\end{axiom}
\begin{axiom}[Some--ax]
\[
\frac{(a:A)}{(\Some(a):\Option(A))}
\]
\[
(I(a)\in \kappa(\Emb(A)))\implies(I(\Some(a))\in\kappa(\Emb(\Option(A))))
\]
\end{axiom}
Completing the predicate-level presentation requires a uniform definition of $\Emb$ and its extension operator $\kappa$; the older class-level equations below are retained with that scope.
\subsubsection{Model of option types}
%Model:\\
\begin{definition}[None--def]
\[
I(\None) \eqdef \emptyclass
\]
\end{definition}
\begin{definition}[Some--def]
\[
I(\Some(a)) \eqdef \{I(a)\}
\]
\end{definition}
\begin{definition}[Option--def](sic!)
\[
\Emb(\Option(A)) \eqdef \cc{b}{(b \sss A)}
\]
\end{definition}
\subsubsection{Soundness of option types}
The following calculations prove the class-level semisingleton model.
They use the empty-set and set-pairing axioms, in addition to comprehension.
Write $O(A)=\{b:b\sss A\}$ for a class $A$. To interpret these as the
proposed typed judgments involving $\Emb$, $\kappa$, $\term$, and
$\type$, those translation operations must identify the option extension
with $O(A)$ and the two constructors with the following class expressions.
Their uniform definitions remain part of the predicate-level proposal;
the class calculation does not itself define them.

\begin{theorem}[None in the semisingleton model]
For every class $A$,
\[
\emptyclass\in O(A).
\]
\end{theorem}
\Proof
The empty class is a set by the empty-set axiom. It is a subclass of
$A$ and has at most one member, since it has none. Thus
$\emptyclass\sss A$, and guarded class comprehension gives
$\emptyclass\in O(A)$. With $I(\None)=\emptyclass$, this proves the
intended None introduction law in this model.
\Done

\begin{theorem}[Some in the semisingleton model]
For every class $A$ and every $a\in A$,
\[
\{a\}\in O(A).
\]
\end{theorem}
\Proof
Membership $a\in A$ implies that $a$ is a set. Pairing therefore makes
$\{a\}$ a set. Its only member is $a$, so it is a semisingleton
subclass of $A$. Comprehension gives $\{a\}\in O(A)$, proving the
intended Some introduction law under $I(\Some(a))=\{I(a)\}$.
Moreover $a\in\{a\}$ makes this code nonempty, so it differs from
None. Equality $\{a\}=\{a'\}$ for set values implies $a=a'$ by
singleton membership. Thus the constructors are disjoint and Some is
injective on the set values represented by $A$.
\Done

\begin{theorem}[The small option carrier and its cases]
In the ordinary classical set model, if $A$ is a set, then $O(A)$
is a set. Every $b\in O(A)$ is either $\emptyclass$ or $\{a\}$
for a unique $a\in A$, and these alternatives are disjoint.
\end{theorem}
\Proof
Every member of $O(A)$ is a set subclass of $A$, so
\[
O(A)=\{b\in\mathcal P(A):b\text{ has at most one member}\}
\]
is a set by powerset and separation. Given $b\in O(A)$, classical
logic distinguishes whether $b$ has a member. If it has none,
extensionality gives $b=\emptyclass$. Otherwise take $a\in b$.
Then $a\in A$, and the at-most-one condition makes every member of
$b$ equal to $a$, so $b=\{a\}$. The earlier constructor properties
give uniqueness of $a$ and disjointness from the empty case.
\Done

Thus this semisingleton carrier is a model of the tagged alternative
``no value or one value'' in the stated classical setting. The case
split is part of the argument: for arbitrary semisingleton subsets in
an intuitionistic setting, decidability of inhabitation cannot be
silently assumed. When $A=\emptyclass$, the carrier is
$O(A)=\{\emptyclass\}$. If $\emptyclass\in A$, Some of that value
is $\{\emptyclass\}$, distinct from None, which is $\emptyclass$.

\begin{theorem}[Dependent option elimination and its computation laws]
Let $A$ be a set and $(E(b))_{b\in O(A)}$ a specified family of sets.
Suppose an element $e_0\in E(\emptyclass)$ and a function $e_1$ on
$A$ are supplied, with $e_1(a)\in E(\{a\})$ for each $a\in A$.
There is a unique function $r$ on $O(A)$ such that
\[
r(b)\in E(b),\qquad
r(\emptyclass)=e_0,\qquad r(\{a\})=e_1(a)\quad(a\in A).
\]
Consequently, reconstructing any section from its values on None and
Some returns that section; this is the semantic eta law for elimination.
\end{theorem}
\Proof
Use ordinary injective set-pair coding to form the graph
\[
r=\{(\emptyclass,e_0)\}
  \cup\{(\{a\},e_1(a)):a\in A\}.
\]
Pairing and replacement make this a set. The two cases have distinct
inputs, and injectivity of singleton formation identifies $a$ whenever
two inputs in the second case agree. Hence the graph is functional.
The case theorem says that its domain is exactly $O(A)$, and the
assumed types of $e_0,e_1(a)$ give $r(b)\in E(b)$.
Its graph directly proves the two computation equations. Any other
section satisfying them agrees with $r$ at the empty input and at
every singleton input; those exhaust the domain, so extensionality
proves uniqueness. For eta, take a given section $q$ and supply
$e_0=q(\emptyclass)$ and $e_1(a)=q(\{a\})$. The section $q$ satisfies
the same equations, so the reconstructed section equals $q$.
\Done

The data $e_0$ and $e_1$ are supplied, rather than chosen from merely
nonempty fibres; this elimination theorem uses no choice axiom.
It proves equality of functions in the set model, not an additional
judgmental conversion rule for the proposed type checker.

\begin{corollary}[Option induction]
Let $A$ be a set and let $Q(b)$ be a predicate on $O(A)$. If
$Q(\emptyclass)$ holds and $Q(\{a\})$ holds for every $a\in A$,
then $Q(b)$ holds for every $b\in O(A)$.
\end{corollary}
\Proof
The case theorem makes each $b$ either the empty set or one of the
singletons covered by the hypotheses. Substitute that equality into
its corresponding hypothesis.
\Done

\pagebreak
\subsection{Implementing ordered pairs}
The ordered-pair section gives the intended constructor, sethood condition, and injectivity property, and compares several candidate models. The type-theoretic presentation here relies on the selected model satisfying that interface.
%\dots\\
\subsubsection{Axioms of ordered pairs}
Use the constructor, sethood, and injectivity interface stated in the ordered-pair section.
\subsubsection{Model of ordered pairs}
Choose one of the pair constructions described there, with its stated domain.
\subsubsection{Soundness of ordered pairs}
This requires the constructor and projection equations for that chosen model.
\pagebreak

\subsection{Implementing list as a medium type}
%of classes[not good]}
Lists of classes are modelled as a family with some initial (and finite) segment of ordinals as a domain.
The type $\ClassList$ is introduced.
\subsubsection{Axioms}
\begin{axiom}[nil]
\[
[] : \ClassList
\]
\end{axiom}
\begin{axiom}[cons]
String ``::'' is the infix notation of a constructor ``cons''.
Parentheses around it make the prefix notation out of infix notation.
\[
(::):\ \Class \to \ClassList \to \ClassList
\]
\end{axiom}
\begin{axiom}[concat]
\[
(\concat):\ \ClassList \to \ClassList \to \ClassList
\]
\end{axiom}
\subsubsection{Model}
%Model:\\
For a medium type, some term interpretations may be proper classes. The model is therefore a predicate on classes, rather than necessarily a set or class containing those interpretations as members.
\begin{definition}[inClassList]
The historical representation proposal for class-list terms is:
%\ClassList\eqdef ? %\cc{(F\family)}{\dom(F)\in \omega}
\[
(F\text{ in ClassList})
\iffdef
(
(F\family)
\land
(\Dom(F)\in\omega)
)
\]%\omega\cup\{\omega\}
\end{definition}
Interpretation of list expressions:
\[
I:\ClassList.\mathsf{Expression}\to\Class.\mathsf{Expression}.
\]
\begin{definition}[nil]
\[
I([]) \eqdef \emptyclass
\]
\end{definition}

\subsubsection{A canonical model with arbitrary class entries}
The preceding predicate needs a fixed meaning of ``family.'' In the
class-axiom chapter, a set value has a single concentrated record;
the semisingleton encoding instead retains an empty record and all
singleton member-records. Thus the latter's cons operation need not
preserve the former family predicate. We now specify a local model
using the canonicalization $C_1$ from the historical family chapter,
with the same ordinary injective set-pair coding throughout this model.
This makes the representation choice explicit while retaining arbitrary
classes as possible entries.

For a class relation $F\subseteq\V\times\V$, write
\[
U_F(i)=\bigcup\{s:(i,s)\in F\},\qquad N(F)=C_1(F).
\]
By the canonicalization theorem, $N$ preserves the represented domain
and every union value, is idempotent, and satisfies
\[
N(F)=\{(i,\emptyclass):i\in\Dom(F)\}
 \cup\{(i,\{q\}):i\in\Dom(F),\ q\in U_F(i)\}.
\]
Define the canonical class-list predicate by
\[
\operatorname{CL}(F)\iff
 F\subseteq\V\times\V\ \land\ N(F)=F\ \land\ \Dom(F)\in\omega.
\]
For such an $F$, its length is the unique natural ordinal
$n=\Dom(F)$. Every relation with finite ordinal domain can be converted
to this model by $N$, preserving its entries and length. In particular
this applies to the historical family proposal, even though the resulting
records can differ. Equal lengths and equal values at every represented
index determine equal canonical lists, by the complete-normal-form
theorem for $C_1$.

\begin{definition}[Class-parameter cons and concatenation]
For any class $H$ and canonical class lists $F,G$, put
\[
\begin{aligned}
\operatorname{Cons}(H,F)
 &=\{(0,\emptyclass)\}\cup\{(0,\{q\}):q\in H\}
   \cup\{(i+1,s):(i,s)\in F\},\\
\operatorname{App}(F,G)
 &=F\cup\{(n+i,s):(i,s)\in G\},\qquad n=\Dom(F).
\end{aligned}
\]
Nil is $\emptyclass$, agreeing with $I([])$ above. The numbers and
their sums in these formulas are natural ordinals.
\end{definition}

These are operations defined by formulas with class parameters.
Every stored coordinate is a set: $q\in H$ makes $q$ a set, hence
$\{q\}$ and its record are sets, even when $H$ is proper. Class
comprehension forms the displayed results; it does not collect a pair
whose coordinate is the proper class $H$ or a proper class list.
No ordinary class function graph on all such arguments is asserted.

\begin{theorem}[Class-list domains and values]
Nil is canonical with domain zero. If $F,G$ are canonical with lengths
$n,m$, respectively, then $\operatorname{Cons}(H,F)$ and
$\operatorname{App}(F,G)$ are canonical, and
\[
\begin{aligned}
\Dom(\operatorname{Cons}(H,F))&=n+1,&
U_{\operatorname{Cons}(H,F)}(0)&=H,\\
&&U_{\operatorname{Cons}(H,F)}(i+1)&=U_F(i)\quad(i<n),\\
\Dom(\operatorname{App}(F,G))&=n+m,&
U_{\operatorname{App}(F,G)}(i)&=U_F(i)\quad(i<n),\\
&&U_{\operatorname{App}(F,G)}(n+j)&=U_G(j)\quad(j<m).
\end{aligned}
\]
\end{theorem}
\Proof
The empty relation has empty domain and is fixed by $N$.
For cons, index zero has exactly the empty record and all singleton
records from $H$. Their union is $H$: each $q\in H$ occurs in
$\{q\}$, and no other element is contributed. For each $i<n$,
the fibre at $i+1$ is exactly the fibre of $F$ at $i$. Successor is
injective on natural ordinals and never zero, so these fibres do not
overlap. They are all nonempty because each contains the empty record.
The domain is therefore $\{0\}\cup\{i+1:i<n\}=n+1$.
Every fibre has precisely the form prescribed by $N$, proving canonicality
as well as the asserted values.

For concatenation, every $k<n+m$ lies uniquely either below $n$ or
has the form $n+j$ with $j<m$. Thus the first part of the union gives
exactly the fibres of $F$ below $n$, and the shifted part gives exactly
the fibres of $G$ at the remaining indices. Every such fibre is
nonempty and canonical, and there are no others. This proves the domain,
canonicality, and both value equations, including the cases $n=0$ or
$m=0$.
\Done

\begin{corollary}[Class-list concatenation identities and associativity]
For canonical class lists $F,G,K$,
\[
\operatorname{App}(\emptyclass,F)=F
 =\operatorname{App}(F,\emptyclass),\qquad
\operatorname{App}(\operatorname{App}(F,G),K)
 =\operatorname{App}(F,\operatorname{App}(G,K)).
\]
\end{corollary}
\Proof
Concatenating with nil leaves the domain and every represented value
unchanged by the theorem, so canonical extensionality gives both
identities. For associativity, write the three lengths as $n,m,p$.
Both outputs have length $(n+m)+p=n+(m+p)$. Splitting the indices
into the consecutive intervals of lengths $n,m,p$, the theorem gives
respectively the values of $F,G,K$ on each side. On the final interval
the two shifts agree because $(n+m)+j=n+(m+j)$. Hence all values and
domains agree, and canonical extensionality proves equality.
\Done

For example, $\operatorname{Cons}(\V,\emptyclass)$ is a one-entry
list whose value is a proper class, represented entirely by set records.
For the empty class as head, the one-entry list is instead
$\{(0,\emptyclass)\}$; it is distinct from nil despite its empty value.
On noncanonical input relations, the corresponding operations first
apply $N$, so an identity law returns $N(F)$ rather than necessarily
the original records $F$. We have proved class-parameter construction
and equality laws, not an initial-algebra property for a collection of
all proper-class-valued lists. The set-valued model below uses these
same canonical fibres and admits a separate set-level recursion theorem.

The newer Obelisk export's entry 3.51, \texttt{DefClassLists}, records
another construction using explicit markers and one-based positions.
The present model uses zero-based positions and ordinary set-pair
records. These are different representations of the intended entries;
neither literal equality of their codes nor interchangeability of their
projection formulas follows. The ordered-pair chapter explains the
separate marker convention. Transport between the encodings must
preserve length and decoded entries explicitly.
%\subsection{Self--contained presentation of dependent types}
%\subsection{Families of definable classes, abstraction and evaluation.}
\subsubsection{Transport to the catalog's finite marker lists}
The newer Obelisk catalog's \texttt{DefClassLists} (3.51), still marked
\texttt{TBI}, proposes a different representation. Its displayed payload
$\{\{\{a\}\},k\}$ is an \emph{unordered} pair: one member is
$\{\{a\}\}$ and the other is the tag $k$. To keep the brace depth
explicit, put
\[
e(a)=\{\{a\}\},\qquad
\operatorname{mark}_2(H,k)
 =\{\{0,k\}\}\cup\{\{e(a),k\}:a\in H\}.
\]
Here $a$ ranges over set members of $H$, while $k$ is a positive finite
ordinal. The subscript records the two singleton operations in $e$.
This is the convention of entry 3.51; the separate tagged-pair entry
\texttt{DefDunaevOrderedPair} (3.76) uses three singleton operations
in its payload. The two formulas are not asserted to be literally equal.
We prove the finite-list interpretation of 3.51, rather than inferring
its correctness from the generated entry's label.

\begin{lemma}[Separation of markers and payloads]
For sets $a,b$ and positive natural ordinals $k,l$, the following hold:
\[
\begin{gathered}
\{0,k\}=\{0,l\}\iff k=l,\\
\{e(a),k\}=\{e(b),l\}\iff (a=b\text{ and }k=l),\\
\{0,k\}\ne\{e(a),l\}.
\end{gathered}
\]
\end{lemma}
\Proof
The set $e(a)$ is a singleton and cannot be an ordinal. Indeed, the
only singleton ordinal is $1=\{0\}$: a nonempty ordinal contains zero,
so its sole member must be zero. Equality $e(a)=1$ would force
$\{a\}=0$, contrary to $a\in\{a\}$. Thus a payload
$\{e(a),k\}$ has exactly one natural-ordinal member, namely $k$,
and exactly one nonordinal member, namely $e(a)$. Equality of two
payloads identifies these members separately. Two applications of
singleton injectivity then give $a=b$. A marker $\{0,k\}$ has
only ordinal members, so cannot equal a payload; its unique nonzero
member is $k$, proving marker injectivity. The reverse implications
follow directly from equality substitution.
\Done

For a canonical class list $F$ of length $n$, define its marker encoding
by a single class comprehension, written as the following union:
\[
\mathcal M(F)=\{\{0,i+1\}:i<n\}
 \cup\{\{e(a),i+1\}:i<n,\ a\in U_F(i)\}.
\]
All members of this class are sets. The expression uses the class
parameter $F$ and its definable values, rather than forming a set or
class whose members are those possibly proper-class values.
For a class $L$ and positive natural $k$, define the decoded entry
\[
D_k(L)=\{a:\{e(a),k\}\in L\}.
\]
As usual comprehension retains set values of $a$; its auxiliary binders
are fresh for the class parameters.

\begin{theorem}[Faithful transport between the two list representations]
For canonical class lists $F,G$, the marker encoding retains both
length and entries. More explicitly, if $\Dom(F)=n$, then
\[
\{k\in\omega\setminus\{0\}:\{0,k\}\in\mathcal M(F)\}
 =\{i+1:i<n\},
\qquad D_{i+1}(\mathcal M(F))=U_F(i)\quad(i<n).
\]
Consequently $\mathcal M(F)=\mathcal M(G)$ if and only if $F=G$.
For arbitrary classes $H,K$ it also satisfies
\[
\begin{aligned}
\mathcal M(\emptyclass)&=\emptyclass,\\
\mathcal M(\operatorname{Cons}(H,\emptyclass))
 &=\operatorname{mark}_2(H,1),\\
\mathcal M(\operatorname{Cons}(H,\operatorname{Cons}(K,\emptyclass)))
 &=\operatorname{mark}_2(H,1)\cup\operatorname{mark}_2(K,2).
\end{aligned}
\]
\end{theorem}
\Proof
By separation of markers and payloads, the marker $\{0,k\}$ occurs
exactly when $k=i+1$ for some $i<n$. The same lemma says that a
payload $\{e(a),i+1\}$ occurs exactly when $i<n$ and $a\in U_F(i)$.
These are the first two assertions by extensionality. In particular an
empty entry still contributes its marker and cannot shorten the list.
The finite tag set is empty exactly when $n=0$, and otherwise has
maximum $n$; hence it determines $n$ uniquely.

If two encodings are equal, their marker sets give equal lengths and
the decoding formulas give equal entries at every represented index.
The $C_1$ normal-form theorem then gives $F=G$. Conversely equal
canonical lists give equal marker comprehensions. Finally nil has no
indices; the one-entry cons list has domain $1$ and value $H$ at zero;
the two-entry cons list has domain $2$ and values $H,K$ at zero and
one. Substituting these domain and value equations into $\mathcal M$
gives the three displayed identities.
\Done

The last two identities give precisely the catalog's singleton-list and
two-entry-list formulas, and $D_1,D_2$ give its stated projections on
their represented indices. They concern a translation of representations,
not equality of a semisingleton-record list with its marker encoding.
The proof allows proper-class entries but only finite natural tags;
it does not verify any unstated kernel conversion rules or an extension
to arbitrary class-valued tags.

\subsection{Implementing list of values having a small type}
Fix a set $A$ interpreting the small type. We use the author's
semisingleton family encoding, rather than the concentrated-fibre
convention used in some other sections. Work with Infinity, powerset,
replacement, union, and ordinary injective pairs of sets. Every value
$h\in A$ is a set; no proper-class entry is being stored in this model.
The displayed type interface below will follow from the construction.

\subsubsection{The encoded list carrier}
For a set $h$, write
\[
\mathcal S(h)=\{s\subseteq h:s\text{ has at most one member}\}.
\]
Thus $s\sss h$ means $s\in\mathcal S(h)$. This is a set by
powerset and separation. It contains $\emptyclass$ and every singleton
$\{q\}$ for $q\in h$, and
\[
\bigcup\mathcal S(h)=h.
\]
Indeed each member of a semisingleton subclass lies in $h$, and each
$q\in h$ occurs in its singleton subclass. In particular
$\mathcal S(\emptyclass)=\{\emptyclass\}$, not the empty family.

For an ordinary function graph $f:n\to A$, where $n\in\omega$, define
\[
L(n,f)=\{(i,s):i\in n,\ s\in\mathcal S(f(i))\}.
\]
Put
\[
\mathcal L_n(A)=\{L(n,f):f:n\to A\},\qquad
\List(A)=\bigcup_{n\in\omega}\mathcal L_n(A).
\]
The graph $f$ is an auxiliary code. The list value is $L(n,f)$,
which generally has several records at each index and is not an
ordinary single-valued graph. Here list evaluation means union of the
fibre, and its represented domain means the indices having a record.

\begin{lemma}[Sethood and recovery of the sequence code]
The collection $\List(A)$ is a set. For every $f:n\to A$,
\[
\Dom(L(n,f))=n,\qquad
\bigcup\{s:(i,s)\in L(n,f)\}=f(i)\quad(i\in n).
\]
Consequently $L(n,f)=L(m,g)$ if and only if $n=m$ and $f=g$.
\end{lemma}
\Proof
For each $i\in n$, the set of records
$\{(i,s):s\in\mathcal S(f(i))\}$ is a set by replacement.
Replacement over $n$ and union make $L(n,f)$ a set. The ordinary
function space $A^n$ is a set, as a separated subclass of
$\mathcal P(n\times A)$. Replacement therefore makes $\mathcal L_n(A)$
a set. Replacement over $\omega$ and union give $\List(A)$.

Every $i\in n$ has the record $(i,\emptyclass)$, and no other index
occurs, proving the domain equation. Injectivity of pair coding makes
the fibre at $i$ precisely $\mathcal S(f(i))$, whose union is $f(i)$.
Equality of two encoded lists therefore gives equal domains, and then
equal values at every index. Function extensionality gives $f=g$.
The reverse implication follows from the defining comprehension.
\Done

The unique $n$ in this lemma is the length $\ell(t)$ of the list
$t=L(n,f)$. An empty entry cannot disappear: its fibre contains the
empty record. This domain information distinguishes the one-entry list
with entry $\emptyclass$ from the empty list.

\subsubsection{Constructors and their computation rules}
\begin{definition}[Nil and the author's cons operation]
Put $[]=\emptyclass$ and retain the definitions
\[
\mathrm{inc\_fst}(r)
 =\{(i\cup\{i\},s):(i,s)\in r\},
\qquad
h::t=\{(\emptyclass,s):s\sss h\}\cup\mathrm{inc\_fst}(t).
\]
On a list's natural-number indices, $i\cup\{i\}=i+1$.
\end{definition}

\begin{theorem}[List constructors and unique decomposition]
One has $[]\in\List(A)$, and $h::t\in\List(A)$ for every
$h\in A$, $t\in\List(A)$. Their lengths satisfy
\[
\ell([])=0,\qquad \ell(h::t)=\ell(t)+1.
\]
Every list is either $[]$ or has the form $h::t$ for unique
$h\in A$ and $t\in\List(A)$, and these cases are disjoint.
Thus the constructors realize the intended interface
\[
[]:\List(A),\qquad (::):A\to\List(A)\to\List(A).
\]
\end{theorem}
\Proof
The unique graph $f:0\to A$ is empty and $L(0,f)=\emptyclass$,
so nil is a list. For $t=L(n,f)$, define the ordinary graph
$g:n+1\to A$ by
\[
g(0)=h,\qquad g(i+1)=f(i)\quad(i\in n).
\]
Every nonzero natural number below $n+1$ has a unique predecessor
below $n$. The record at index zero comes from the first term of cons,
and the records at index $i+1$ are precisely those in the fibre of
$t$ at $i$. Hence the defining equations give $h::t=L(n+1,g)$.
This proves closure and the length formulas. In particular every cons
has $(0,\emptyclass)$ as a record, so it cannot equal nil.

Conversely, if $L(n,f)$ is not nil, its domain is not zero. Write
$n=m+1$, put $h=f(0)$, and let $f_-(i)=f(i+1)$ for $i\in m$.
The same fibre calculation gives
$L(n,f)=h::L(m,f_-)$. Recovery of the sequence code proves uniqueness:
the head is the value at zero, and all tail entries are the subsequent
values in their order. This also proves that $h::t=h'::t'$ implies
$h=h'$ and $t=t'$.
\Done

For example, if $\emptyclass\in A$, then
$\emptyclass::[]=\{(0,\emptyclass)\}$ has length one, although its
only evaluated entry is empty. No out-of-domain default is used to
recover its length. If $A=\emptyclass$, there are no positive-length
sequence graphs into $A$, so $\List(A)=\{[]\}$.

\subsubsection{Induction and recursion on length}
\begin{theorem}[List induction and recursion]
A predicate that holds at nil and is preserved by cons holds at every
member of $\List(A)$. Given a set $X$, an element $z\in X$, and a
function $c:A\times X\to X$, there is a unique function
$r:\List(A)\to X$ satisfying
\[
r([])=z,\qquad r(h::t)=c(h,r(t)).
\]
\end{theorem}
\Proof
For induction, induct on the natural-number length. The zero-length
list is nil, and every list of length $n+1$ uniquely decomposes as
$h::t$ with $\ell(t)=n$. The corresponding induction premise and
cons step prove the assertion at that length.

For recursion, define functions $r_n:\mathcal L_n(A)\to X$.
On $\mathcal L_0(A)=\{[]\}$ set $r_0([])=z$. Given $r_n$, the
unique decomposition lets us define
\[
r_{n+1}(h::t)=c(h,r_n(t))\qquad(t\in\mathcal L_n(A)).
\]
Replacement collects this graph, and natural-number recursion produces
the sequence $(r_n)$. Distinct length classes are disjoint by code
recovery, so $r=\bigcup_{n\in\omega}r_n$ is a function with domain
$\List(A)$ and values in $X$. Replacement and union make its graph
a set. The stage equations give the required computation rules.
Any other solution agrees on length zero, and agreement on length $n$
forces agreement on length $n+1$ by unique decomposition. Induction
therefore gives equality everywhere, proving uniqueness.
\Done

This is the usual initial-algebra property for the tagged sum functor
$X\mapsto 1+(A\times X)$ on sets: the nil and cons branches specify
an algebra, and the two recursion equations specify its unique morphism
from the encoded list algebra. The disjoint alternatives are supplied
by the tags in that functor, while the represented list domain already
distinguishes nil from cons in the concrete encoding. This result does
not assert an initial-algebra property for the class-parameter model
in the preceding subsection.

\subsubsection{Concatenation and maps of small lists}
For each $u\in\List(A)$, apply list recursion with target $\List(A)$,
nil value $u$, and step $(h,v)\mapsto h::v$. Write its value at $t$
as $t\mathbin{+\!\!+}u$. Thus
\[
[]\mathbin{+\!\!+}u=u,\qquad
(h::t)\mathbin{+\!\!+}u=h::(t\mathbin{+\!\!+}u).
\]
Uniqueness specifies the operation for every pair $(t,u)$, and replacement
on $\List(A)\times\List(A)$ makes it a set function. This realizes
the earlier concatenation interface for small lists in this encoding.

\begin{proposition}[Concatenation laws]
For $t,u,v\in\List(A)$,
\[
\begin{gathered}
t\mathbin{+\!\!+}[]=t,\qquad
(t\mathbin{+\!\!+}u)\mathbin{+\!\!+}v
 =t\mathbin{+\!\!+}(u\mathbin{+\!\!+}v),\\
\ell(t\mathbin{+\!\!+}u)=\ell(t)+\ell(u),\\
t\mathbin{+\!\!+}u=t\mathbin{+\!\!+}v\quad\Longrightarrow\quad u=v.
\end{gathered}
\]
Together with the defining nil equation, nil is a two-sided identity.
\end{proposition}
\Proof
For the right identity, the function $t\mapsto t\mathbin{+\!\!+}[]$
and the identity on $\List(A)$ both take nil to nil and satisfy the
same cons equation. Recursion uniqueness makes them equal.

Fix $u,v$. As functions of $t$, both sides of the associativity equation
take nil to $u\mathbin{+\!\!+}v$. Applying either function to $h::t$
gives $h$ consed onto its value at $t$, by the two defining equations.
Recursion uniqueness again proves equality.

For length, induct on $t$. At nil the equation is
$\ell(u)=0+\ell(u)$. At $h::t$, concatenation and the length formula
give $\ell(t\mathbin{+\!\!+}u)+1$; the induction hypothesis makes
this $\ell(t)+\ell(u)+1=(\ell(t)+1)+\ell(u)$, as required by
natural-number arithmetic. Finally, induct on the common left list
for cancellation. At nil the hypothesis is already $u=v$.
At $h::t$, constructor injectivity cancels $h$ from
$h::(t\mathbin{+\!\!+}u)=h::(t\mathbin{+\!\!+}v)$, and the
induction hypothesis gives $u=v$.
\Done

For a set function $f:A\to B$, list recursion with target $\List(B)$,
nil value $[]$, and step $(h,v)\mapsto f(h)::v$ defines
$\operatorname{map}_f:\List(A)\to\List(B)$, with
\[
\operatorname{map}_f([])=[],\qquad
\operatorname{map}_f(h::t)=f(h)::\operatorname{map}_f(t).
\]
These equations use the same semisingleton list encoding on both sides;
the operation replaces list entries, not individual records in a fibre.

\begin{proposition}[Mapping identities, composition, and concatenation]
For set functions $f:A\to B$, $g:B\to C$,
\[
\begin{gathered}
\operatorname{map}_{1_A}=1_{\List(A)},\qquad
\operatorname{map}_{g\circ f}=\operatorname{map}_g\circ\operatorname{map}_f,\\
\operatorname{map}_f(t\mathbin{+\!\!+}u)
 =\operatorname{map}_f(t)\mathbin{+\!\!+}\operatorname{map}_f(u).
\end{gathered}
\]
\end{proposition}
\Proof
For the identity equation, both functions send nil to nil and send
$h::t$ to $h$ consed onto their value at $t$. For composition, both
send nil to nil and send $h::t$ to $g(f(h))$ consed onto their value
at $t$. Each comparison is therefore an application of recursion
uniqueness with the indicated target set and step operation.

For the last equation fix $u$. As functions of $t$, both sides take
nil to $\operatorname{map}_f(u)$. The cons and mapping equations make
both send $h::t$ to $f(h)$ consed onto their value at $t$. Recursion
uniqueness with that nil value proves the claimed compatibility.
\Done

All constructions in this subsection have set-valued entries and set
carriers. The medium-type concatenation laws above have their separate
class-parameter proof; the present recursion and mapping arguments do
not assert a recursion principle for arbitrary class-valued entries.
\pagebreak


\subsection{Implementing functions}
This subsection records class operations and compares two canonical family
encodings. A relation is a class of encoded set pairs. A family may use
several entries in a fiber to represent a proper-class value, so relational
composition and family evaluation must not be interchanged without checking
the representation. The historical conversion laws below explicitly return
to semisingleton fibers and total union evaluation.

%\subsection{Relational basics}
\subsubsection{Family basics}
\iffalse %OUTDATED
%нету "неопределённого" значения
\begin{definition}[value]
\[
F(a)\eqdef\union\cc{s}{(a,s)\in F}
\]
\end{definition}
\fi %OUTDATED
%2
\iffalse %it doesn't matter whether it is set or class, bad!
%Also functions are now relations too!
\begin{definition}[relational domain]
\[
\dom(R)\eqdef\cc{a}{\exists b,\ (a,b)\in R}
\]
\end{definition}
\begin{definition}[functional Domain]%WRONG: \Dom(F)\eqdef\dom(F)
\[
\dom(R)\eqdef\cc{a}{\exists (b\set),\ (a,b)\in R}
\]
\end{definition}
\fi %OUTDATED
%3
\iffalse %OUTDATED %"relational"
\begin{definition}[relational range]
\[
\ran(R)\eqdef\cc{b}{\exists a,\ (a,b)\in R}
\]
\end{definition}
\fi %OUTDATED
\iffalse bad
\begin{definition}[functional Range]
\[
\Ran(F)\eqdef\cc{F(a)}{a\in\Dom(F)}
\]
\end{definition}
\fi
%4
\begin{definition}[restriction]%OK
\[
F\restriction A\eqdef\cc{(a,s)\in F}{a\in A}
\]
\end{definition}
%5
\begin{definition}[image]%relational image]
\[
\image_R(A)\eqdef\ran(R\restriction A)
\]
\end{definition}
\iffalse
\begin{definition}[functional Image]
\[
\Image_F(A)\eqdef\Ran(F\restriction A)
\]
\end{definition}
\fi
%6
\iffalse %OUTDATED (it doesn't work with families)
\begin{definition}[relational composition]
\[
R \circ S \eqdef \cc{(x,z)}{\exists y,\ (x,y)\in S \land (y,z)\in R}
\]
\end{definition}
\fi %OUTDATED
\begin{definition}[composition]%functional composition]
\[
F \circ G \eqdef \cc{(x,z)}{x \in \Dom(G) \land (G(x),z)\in F}
\]% z may be either a singleton subset of a proper value, or a value.
\end{definition}
%7
\iffalse %OUTDATED (probably already defined [as class])
\begin{definition}[class of all sets]
\[
\V\eqdef\cc{x}{\top}%x\any
\]
\end{definition}
\fi %OUTDATED
\begin{definition}[membership class]
\[
\E\eqdef\cc{(x,y)}{x\in y}
\]
\end{definition}
\begin{definition}[identity class]
\iffalse
\[
\I\eqdef\cc{(x,x)}{x\any}
\]
\fi
\[
\I\eqdef\ccc{x}{(x,x)}{\top}
\]
\end{definition}
%8
\iffalse %OUTDATED
\begin{definition}[relational transposition]
\[
R^\top\eqdef\cc{(y,x)}{(x,y)\in R}
\]
\end{definition}
\fi %OUTDATED
\begin{notation}[containment relation]
\[
\invE \eqdef \E^\top
\]
\end{notation}
\begin{definition}[transposition]%functional transposition]
\[
F^\intercal\eqdef\cc{(F(x),x)}{x\in \dom(F)}
\]
\end{definition}

\subsubsection{Family}
A family combines its domain with the chosen fiber representation. The preceding family-evaluation section gives the two encodings compared here.
\iffalse %OUTDATED
\begin{definition}[family]
Def is ommited!
\end{definition}

\begin{definition}[domain]%WRONG it is Dom.
\[
\dom(F)\eqdef \cc{x}{\exists b,\ (x,b)\in F}
\]
\end{definition}

\begin{definition}[value]% again, no "undefined" case
\[
F(x) \eqdef \bigcup\cc{c}{((x,c)\in F)}
\]
\iffalse
\[
F(x) \eqdef \cc{w}{\exists b, (w\in b) \land ((x,b)\in F)}
\]

% It may have sense to generalize value, but (x,b)\in F \dots
% what can be done if b is a proper class?
% if (x,b) = ({0}\times x)\cup({1}\times b), then
% \dots. wait! Family is already a generalisation for the same purpose!
(small idea:
\[
?\eqdef \bigcup^Q\limits_{Q\subclass \power Q} Q.
\]
)
\dots
\fi
\end{definition}
\fi %OUTDATED

\subsubsection{Lambda--abstraction}
%BADRecall:
%\[
%\cc{(x,t)}{(x\in A)}
%= \cc{p}{\exists x, (p = (x,t)) \land (x\in A)}
%\]
\begin{baddefinition*}[set--valued abstraction]
This graph abstraction represents set values. It can be used when $t$ is set-valued on $A$; extending it to proper-class values requires another encoding:
\[
(x\in A \mapsto t)
\eqbaddef
\ccc{x}{(x,t)}{(x\in A)}
\clabeltagd{def}{setvaluedabstraction}
\]
(Here $x$ is a bound variable on the both sides of the equality.)
%a binder
\end{baddefinition*}

In both following abstractions, $x$ is not free in the fixed domain $A$,
and the auxiliary variable $b$ is distinct from $x$ and fresh for $A,t$.
Rename auxiliary variables before substitution whenever needed.

\begin{definition}[OUTDATED class--valued abstraction]
\[
(x\in A \mapsto t)
\eqdef
\ccc{x,b}{(x,b)}{(x\in A)\land (b \sss t)}
\clabeltagd{def}{classvaluedabstraction-old}
\]
\end{definition}
In the singleton-based variant, a set value is stored directly and a proper-class value is represented by singleton subclasses:
\begin{definition}[class--valued abstraction]
\[
(x\in A \mapsto t)
\eqdef
\ccc{x,b}{(x,b)}{
(x\in A)\land
(
%\lif\
((t\set)\land(b=t))
\lor
%\ \lthen\
((t\proper)\land(b\singleton)\land(b\subset t))
)
}
\clabeltagd{def}{classvaluedabstraction}
\]
\end{definition}

%\mathclap{#1}
\iffalse
%\ |^x\ \ \smash{\stackrel{x}{|}}

\newlength{\heightofhw}
\settoheight{\heightofhw}{$
\vphantom{
\left\{\,
\middle|
2\  \middle|^{1} |^1 \middle|\ 3\,
\right\}
}$
}

%\fontsize{0.04\textheight}{0.04\textheight}
%\the\heightofhw
%\measuredelimiter{
\[
\left(
jjj
\left\{\,
\middle|
2\  \middle|^{1} |^1 \middle|\ 3\,
\right\}
\right)
\]
%}
\fi

% \def\delimsize{\middle}
%\heightof{ }
% \rule{width}{height} to match exact height
% \scalebox{0.8}{}
% q:BTerm
% q:= (x.t)
% lam q := cc (x. exists x, (p = (x,t)) \land (x\in A) )
\subsubsection{Alpha--equivalence}

Renaming the bound variable in an abstraction preserves its meaning, provided the new variable is fresh and the renaming avoids capture. This follows from the same renaming property for class comprehension.

\subsubsection{Beta--conversion}
\begin{definition}[conditional operator]
For any formula $\phi$ and for any class terms $t_1$ and $t_2$, the conditional operator is defined in the following way:
\[
(\lif\ \phi\ \lthen\ t_1\ \lelse\ t_2\ \lfi)
\eqdef
\cc{w}{((\phi)\land(w\in t_1))\lor((\neg\phi)\land(w\in t_2))}
\]
\underline{Side conditions:}\\
1) $w$ is not free in $\phi$,\\
2) $w$ is not free in $t_1$,\\
3) $w$ is not free in $t_2$.
\end{definition}


\begin{theorem}[OUTDATED $\beta$--conversion]
This theorem assumes the outdated implementation:
\[
F(x) \eqdef \bigcup\cc{c}{((x,c)\in F)}
\]

All auxiliary fibre variables below are chosen fresh for $A,t,a,F$;
the inner abstraction uses a different variable $c$ from the outer fibre
variable $b$. Let $A$ be a class term with $x$ not free in $A$, let
$a$ be a set-valued argument term, and let substitution of $a$ for $x$
avoid capture. The body value $t[a/x]$ may be a proper class. For the
semisingleton abstraction,
\[
(x\in A \mapsto t)(a) = (\lif\ (a\in A)\ \lthen\ t[a/x]\ \lelse\ \emptyclass\ \lfi)
\clabeltagd{thm}{fam--beta}
\]
\Proof
\[
(x\in A \mapsto t)(a)
\]
\[
=
\bigcup\cc{b}{((a,b)\in (x\in A \mapsto t))}
\]
\[
=
\bigcup\cc{b}{((a,b)\in \ccc{x,c}{(x,c)}{(x\in A)\land (c \sss t)})}
\]
\[
=
\bigcup\cc{b}{(a\in A)\land (b \sss t[a/x])}
\]
\[
=
\bigcup (\lif\ (a\in A)\ \lthen\ \cc{b}{(b \sss t[a/x])}\ \lelse\  \emptyclass\ \lfi)
\]
\[
=
(\lif\ (a\in A)\ \lthen\ \bigcup\cc{b}{(b \sss t[a/x])}\ \lelse\  \bigcup\emptyclass\ \lfi)
\]
\[
=
(\lif\ (a\in A)\ \lthen\ t[a/x]\ \lelse\ \emptyclass\ \lfi)
\]
\Done
\end{theorem}
The set-argument hypothesis is the scope in which the chosen set-pair
encoding justifies identifying graph coordinates in this calculation.
It does not restrict the represented body values to sets. An extension
of the equation to proper-class arguments would require a specified
pair encoding and evaluation law for those arguments. The empty result
outside $A$ belongs to this local unguarded union convention.

\subsubsection{Eta--conversion}
\begin{theorem}[OUTDATED $\eta$--conversion]
Assume that $F$ uses canonical semisingleton fibers on its represented
domain $A$, and in this local notation write $\dom(F)=A$. Choose the
abstraction variable $x$ fresh for the class expressions $F$ and $A$;
choose the fibre variables $b,c$ distinct from $x$ and fresh for those
expressions. Thus $F$ remains a fixed family when the binder is introduced.
The abstraction and evaluation use the same local semisingleton and
unguarded union conventions as in the preceding theorem. Then
\[
(x\in \dom(F) \mapsto F(x))=F.
\clabeltagd{thm}{fam--eta}
\]
\Proof
\[
(x\in \dom(F) \mapsto F(x))
\]
\[
=\ccc{x,b}{(x,b)}{(x\in A)\land (b \sss F(x))}
\]
\[
=\ccc{x,b}{(x,b)}{\left(x\in A\middle)\land \middle(b \sss \bigcup\cc{c}{((x,c)\in F)}\right)}
\]
\[
=\ccc{x,b}{(x,b)}{\left(x\in A\middle)\land \middle(b \in \cc{c}{((x,c)\in F)}\right)}
\]
\[
=\ccc{x,b}{(x,b)}{((x,b)\in F)}
\]
\[
=F
\]
\Done
\end{theorem}

%\section{Class theory}
%\subsection{Basic class--theoretic operations}
\subsubsection{Relations and their induced operators}
Fix set carriers $X,Y$ and a relation $R\subseteq X\times Y$.
Write $R[A]=\{y\in Y:\exists x\in A,\ (x,y)\in R\}$ and
$R^\top[B]=\{x\in X:\exists y\in B,\ (x,y)\in R\}$.
The direct-image operator $R^{\rightarrow}:\power X\to\power Y$
and inverse-image operator $R^{\leftarrow}:\power Y\to\power X$
are the functions sending a subset to these respective images.
They are functions even when $R$ is not functional: each subset has a
uniquely specified image, though an individual point may have many outputs.
The identity relation $I_X=\{(x,x):x\in X\}$ is a subset of
$X\times X$, since both coordinates of each record belong to $X$.
Restriction $R|_A$ retains exactly the records whose first coordinate
belongs to $A$.

\begin{theorem}[Images preserve unions]
For $\mathcal M\subseteq\power X$ and
$\mathcal N\subseteq\power Y$,
\[
R[\bigcup\mathcal M]=\bigcup R^{\rightarrow}[\mathcal M],
\qquad
R^\top[\bigcup\mathcal N]=\bigcup R^{\leftarrow}[\mathcal N].
\]
Neither functionality nor injectivity of $R$ is required.
\end{theorem}
\Proof
Membership in the first left side means that some $A\in\mathcal M$
contains an $x$ with $(x,y)\in R$. Equivalently, $y\in R[A]$ for
some $A\in\mathcal M$, which is membership in the right side.
Apply the same argument to $R^\top$ for the second equality.
The empty indexing families give empty images on both sides.
\Done

The fixed set carriers ensure that all image values used as members of
the right-hand families are sets. For arbitrary class relations an image
can be proper, and a class graph cannot store that value as a member.
An extension to that setting must use the family's class-value encoding
or retain the quantified membership formulation instead.

\begin{theorem}[Functional images of countable sets]
If $M$ admits an injection $j:M\to\omega$ and $g$ is a set-valued
class function, then $g[M]$ is countable, where only arguments in the
domain of $g$ contribute to the graph image.
\end{theorem}
\Proof
First $j[M]\subseteq\omega$ is a set by separation, and the inverse
class function $j^{-1}$ on $j[M]$ has image $M$. Replacement therefore
makes $M$ a set and then makes $g[M]$ a set. For each $y\in g[M]$, let $k(y)$
be the least $j(x)$ among $x\in M\cap\dom(g)$ with $g(x)=y$. Such a natural
number exists by the definition of image and well-ordering of $\omega$.
If $k(y)=k(z)$, injectivity of $j$ identifies the two witnessing inputs,
and functionality of $g$ gives $y=z$. Thus $k:g[M]\to\omega$ is
injective. The empty image is included. Least witnesses avoid any appeal
to a choice axiom.
\Done

The union--powerclass comparison
$A\subseteq\power B\iff\bigcup A\subseteq B$ was proved in the
union and intersection section. It is a statement about arbitrary classes;
it is distinct from the set-carrier restrictions needed to form the
image operators here.

%%%%%%%%%%%%%%%%%%%%%%%%% END

%Implementing f
%\subsubsection{Axioms of $\Sigma$--types}
%\begin{definition}(Union)
%\[
%\union A\eqdef\cc{x}{\exists a\in A: x\in a}
%\]
%\end{definition}
%1

%9
%12
%%14
%\begin{definition}(Sigma--class)
%\[
%\]
%\end{definition}
\pagebreak
\subsection{Implementing $\Sigma$--types}
\subsubsection{Axioms of $\Sigma$--types}
The formation rule for a dependent sum is:
\begin{prooftree}
\AxiomC{
$(A:\Type)$\quad\quad$(x:A\vdash B(x):\Type)$
}
\RightLabel{$\Sigma$--form}
\UnaryInfC{$
\vdash \Sigma_{(x:A)}B(x):\Type
$}
\end{prooftree}

\subsubsection{Model of $\Sigma$--types}
Throughout this subsection, $B$ is a family. The dependent sum consists of pairs of an index in its domain and a member of the corresponding value. The domain guard is explicit so that an out-of-domain default contributes no extra pairs.
\iffalse
\begin{definition}[dependent sum]
\iffalse %OUTDATED: functional transposition vs relational transposition
\[
\Sigma B \eqdef \mbox{E}^{\top} \circ B
\]
\fi
\[
\Sigma B \eqdef \cc{(a,b)}{(a\in\Dom(B))\land (b\in B(a)) }
\]
\end{definition}
\fi
\begin{definition}[dependent sum of a family]
\[
\Sigma B\eqdef\cc{(x,y)}{(x\in\Dom(B))\land(y\in B(x))}
\]
\end{definition}
%10,11
\begin{definition}[ordered pair projections]
\[
\pr_1 \eqdef \ccc{a,b}{((a,b),a)}{\ a\any,b\any}
\]
\[
\pr_2 \eqdef \ccc{a,b}{((a,b),b)}{\ a\any,b\any}
\]
\end{definition}




\subsubsection{Soundness of $\Sigma$--types}
\begin{theorem}[Soundness of $\Sigma B$ Formation, Introduction, Projections, and $\beta$/$\eta$ Rules]
Let $B$ be a family of classes over domain $A = \Dom(B)$, and let
\[
\Sigma B \eqdef \cc{p}{\exists (x \in \Dom(B)),\ \exists (y \in B(x)),\ (p = (x, y))}.
\]
Then:
\begin{enumerate}
\item \textbf{Introduction (Pairing):} If $(x \in \Dom(B))$ and $(y \in B(x))$, then $((x, y) \in \Sigma B)$.
\item \textbf{Elimination (Projections):} For every $(p \in \Sigma B)$, writing $p = (x, y)$ with $(x \in \Dom(B))$ and $(y \in B(x))$, ordered-pair injectivity yields $(\pr_1(p) = x \in \Dom(B))$ and $(\pr_2(p) = y \in B(\pr_1(p)))$.
\item \textbf{$\beta$- and $\eta$-Computation:} $(\pr_1((x, y)) = x)$, $(\pr_2((x, y)) = y)$, and $(( \pr_1(p), \pr_2(p) ) = p)$ for all $(p \in \Sigma B)$.
\end{enumerate}
The newer catalog records the dependent-sum membership interface as \texttt{DefSigmaModel} (3.38). The projection and computation claims additionally use the selected pairing model; that entry alone is not a certificate for every rule displayed here.
\end{theorem}

\subsubsection{Rules for small $\Sigma$--types}
For these rules, $\Un$ denotes a chosen universe closed under the dependent
sum construction. That closure is an assumption on this local universe.
\underline{Formation of terms of $\Sigma$--type:}\\
\[
\Sigma : \Pi((A:\Un)\mapsto((A\to \Un)\to \Un))
\]
\underline{Introduction of terms of $\Sigma$--type:}\\
\[
\textrm{pair} : \Pi((A:\Un)\mapsto\Pi((B:A\to \Un)\mapsto\Pi((x:A)\mapsto(B\ x\to\Sigma\ A\ B))))
\]
\underline{Elimination of terms of $\Sigma$--type:}\\
\[
\textrm{pr}_1 : \Pi((A:\Un)\mapsto\Pi((B:A\to \Un)\mapsto ((\Sigma\ A\ B) \to A)))
\]
\[
\textrm{pr}_2 : \Pi((A:\Un)\mapsto\Pi((B:A\to \Un)\mapsto \Pi((p:\Sigma\ A\ B)\mapsto B\ (\pr_1\ p))))
\]
\underline{Computation of terms of $\Sigma$--type:}\\
\[
\pr_1\ (\textrm{pair}\ A\ B\ x\ y) \equiv x
\]
\[
\pr_2\ (\textrm{pair}\ A\ B\ x\ y) \equiv y
\]
\[
\textrm{pair}\ A\ B\ (\pr_1\ p)\ (\pr_2\ p) \equiv p
\]
\subsubsection{The membership expansion of a family}
For an ordinary set-valued class function $B$, the dependent sum has
records $(x,y)$ with $x\in\dom(B)$ and $y\in B(x)$. Hence
\[
\Sigma B=\invE\circ B.
\]
Indeed, a record in the composite supplies a set $z$ with $(x,z)\in B$
and $y\in z$; functionality identifies $z=B(x)$, and the converse
uses that same witness. In particular, for the identity class function
$I=\{(x,x):x\in\V\}$, this gives $\Sigma I=\invE$.
These equalities concern ordinary set-valued graphs; applying the same
notation to a proper-class-valued encoded family requires its separately
specified expansion convention.
\pagebreak
\subsection{Implementing $\Pi$--types}
\begin{notation}[Fm]
The terms for propositions are called \textit{formulas}.
\[
\Fm\equivdef\Prop\Term
\]
\end{notation}
\begin{notation}[typecong]
There are no conglomerate--level terms except conglomerate--level types and virtualisations of them.
That is the reason why ``typecong'' is a keyword, that makes a judgment out of class--level term and conglomerate--level type.
There are exactly two variants of judgements: it is either a typing $(a:A)$ or declaring a conglomerate--level type.
\end{notation}
\begin{notation}[conglomerate--level types]
Any formula (e.g. $\mquote{(x \set)}:\Prop\Term$), which is constructed out of a unary predicate symbol is semantically written like this:
\[
(x \set),
\]
while a typing judgement and variable's binding are syntactically written like this:
\[
(x:\Set).
\]
The former is the valuation of the latter. To be precise, specifically, it holds under one special assignment which maps all letters to corresponding variables of the current context.
\end{notation}
%\[
%\forall (B\family),
%\]
\subsubsection{Axioms of $\Pi$--types}
\underline{Formation of terms of $\Pi$--type:}\\
Class--level type theory is predicative.\\
In class--level type theory, the formation rule for a $\Pi$--type is
\begin{prooftree}
\AxiomC{
$(B \family)$
}
\RightLabel{$\Pi$--form}
\UnaryInfC{$
(\Pi B)\text{ typecong}
$}
\end{prooftree}
In particular,
\begin{prooftree}
\AxiomC{
$(B \text{ class--function})$
}
\RightLabel{$\Pi$--form}
\UnaryInfC{$
(\Pi B)\text{ typecong}% sic! it's not a mistake!
$}
\end{prooftree}
\begin{prooftree}
\AxiomC{
$(B \text{ family})\quad\quad(\Dom(B)\set)$
}
\RightLabel{$\Pi$--form}
\UnaryInfC{$
((\Pi B)\text{ class})$%% COOL!
}
\end{prooftree}%In particular,
\begin{prooftree}
\AxiomC{
$(B \text{ function})$
}
\RightLabel{$\Pi$--form}
\UnaryInfC{$
((\Pi B)\text{ set})
$}
\end{prooftree}
In set--level type theory, the formation rules for a $\Pi$--type are
\begin{prooftree}
\AxiomC{
$\Wau \vdash (A:\text{Class})$\quad\quad$\Wau, (x:A)\vdash (B(x):\text{Class})$
}
\RightLabel{$\Pi$--form(large)}
\UnaryInfC{$
\Wau \vdash \Pi_{(x:A)}B(x)\text{ typecong}%:\text{Class}
$}
\end{prooftree}
\begin{prooftree}
\AxiomC{
$\Wau \vdash (A:\text{Set})$\quad\quad$\Wau, [(x:A)]\vdash (B(x):\text{Set})$
}
\RightLabel{$\Pi$--form(small)}
\UnaryInfC{$
\Wau \vdash (\Pi_{(x:A)}B(x)):\text{Set}
$}
\end{prooftree}
For the tiny formation rule, assume additionally that the chosen universe
$\Un$ is closed under dependent products: if $A\in\Un$ and a specified
family has $B(a)\in\Un$ for all $a\in A$, then its dependent product
belongs to $\Un$. The dependent-sum closure assumed earlier does not by
itself supply this requirement.
\begin{prooftree}
\AxiomC{
$\Wau \vdash (A:\Un)$\quad\quad$\Wau, [(x:A)]\vdash (B(x):\Un)$
}
\RightLabel{$\Pi$--form(tiny)}
\UnaryInfC{$
\Wau \vdash \Pi_{(x:A)}B(x):\Un
$}
\end{prooftree}
The following implementation proposal omits general subtyping and instead permits the two-sort inclusion from Set to Class. This is a local design choice for these rules.
\begin{prooftree}
\AxiomC{
$(A:\text{Set})$
}
\RightLabel{(Set is a subsort of Class)}
\UnaryInfC{$
(A:\text{Class})
$}
\end{prooftree}
In the contextual rules below, $x$ is fresh for $\Gamma$, all displayed
contexts and types are well formed, and substitutions avoid capture.
The bound name in the eta rule is chosen fresh for $f$ as well as for
$\Gamma$. These conventions permit alpha-renaming of local binders.
\underline{Introduction of terms of $\Pi$--type:}\\
\begin{prooftree}
\AxiomC{
$\Gamma, x:A\vdash b:B$
}
\RightLabel{$\Pi$-intro}
\UnaryInfC{
$\Gamma \vdash (\lambda ((x:A)\mapsto b)):\Pi_{(x:A)}B$
%$\vdash \Pi_{(x:A)}B(x):\Type$
}
\end{prooftree}
\underline{Elimination of terms of $\Pi$--type:}\\
\begin{prooftree}
\AxiomC{
$\Gamma \vdash(f:\Pi_{(x:A)}B)$\quad\quad$\Gamma\vdash(a:A)$
}
\RightLabel{$\Pi$-elim}
\UnaryInfC{
$\Gamma \vdash f(a):B[a/x]$
}
\end{prooftree}
\underline{Computation of terms of $\Pi$--type:}\\
\begin{prooftree}
\AxiomC{
$\Gamma, x:A\vdash b:B$
\quad\quad
$\Gamma\vdash(a:A)$
}
\RightLabel{$\Pi$-comp}
\UnaryInfC{
$\Gamma \vdash (\lambda ((x:A)\mapsto b))(a) \equiv b[a/x] : B[a/x]$
}
\end{prooftree}
\underline{Uniqueness of terms of $\Pi$--type:}\\
\begin{prooftree}
\AxiomC{
$\Gamma \vdash f:\Pi_{(x:A)}B$
}
\RightLabel{$\Pi$-uniq}
\UnaryInfC{
$\Gamma \vdash f\equiv (\lambda ((x:A)\mapsto f(x))):\Pi_{(x:A)}B$
}
\end{prooftree}

\subsubsection{The first model of $\Pi$--types as predicates on classes}
\begin{remark}[a section needs more than a common domain]
The newer Obelisk record \texttt{DefPiPredicateModel} (3.39) displays an
implication from ``$B$ and $f$ are families with the same domain'' to
$\forall x\in\Dom(B),\ f(x)\in B(x)$. The pointwise membership condition
must instead be part of the section predicate or a typing hypothesis.
For example, take $B=f=\{(0,\emptyclass)\}$. Both domains are $\{0\}$,
but $f(0)=B(0)=\emptyclass$, so the claimed membership would be
$\emptyclass\in\emptyclass$. The predicate defined below retains the
necessary condition. The generated record's verification label alone does
not resolve this discrepancy in its displayed statement.
\end{remark}

Let $B$ be a family term.\\
So $B\eqdef\mantiquoteleft \pmb{B} \mantiquoteright^{\Family}_v$ is a family of classes.\\
%Let $B$ be a family of classes.\\

For this model the judgement ``$\pmb{(f:\Pi B)}$'' is interpreted as a proposition ``$(f:\Pi B)$''.
\begin{definition}[judgement as a proposition]
\[
\mantiquoteleft
\pmb{(f:\Pi B)}
\mantiquoteright^{\Prop}_v
\iffdef
(
\mantiquoteleft
\pmb{f}
\mantiquoteright^{\Class}_v
:\Pi
\mantiquoteleft
\pmb{B}
\mantiquoteright^{\Class}_v
),
\]
\end{definition}
where a predicate symbol is defined the following way:
\begin{definition}[class--function $f$ respects a family $B$]
\[
(f:\Pi B)
\iffdef
\]
\[
\iff
(
(
f\text{ class--function})
\land
\]
\[
\land
(\dom(f)=\Dom(B))
\land
\]
\[
\land
(\forall (x\in\dom(f)), (f(x)\in B(x)) )
)
\]
\end{definition}
This predicate presentation allows the domain of $B$ to be a proper class. Any function with that domain must then itself be proper, and so cannot be a member of a class. Keeping $(f:\Pi B)$ as a predicate avoids imposing an unintended sethood condition on $f$.


\subsubsection{The second model of $\Pi$--types as sets}
%Assume that $B$ is a function. NO
The family terms are valuated as families. %(todo: prove it)
Let family $B$ be the valuation of a family term $\pmb{B}$.
Note, that
\[
\mantiquoteleft
\pmb{\Dom(B)}
\mantiquoteright^{\Class}_v
=
\Dom(
\mantiquoteleft
\pmb{B}
\mantiquoteright^{\Family}_v
).
\]
Let $B\eqdef\mantiquoteleft \pmb{B} \mantiquoteright^{\Family}_v$.
For a class of set-valued sections it is enough to use class comprehension. For the displayed interpretation specifically as a set, assume in addition that $B$ is a set-indexed family of sets, so that its dependent sum and the relevant function space are sets. Then
\[
\mantiquoteleft
\pmb{\Pi B}
\mantiquoteright^{\Set}_v
\eqdef
\cc{f}{(f:\Pi B)}
\]


\subsubsection{Further class-valued models of $\Pi$--types}
The following variants collect set-sized sections. On a proper-class domain,
they need not represent all instances of the predicate model. The size
terminology in this proposal is local: it distinguishes syntax-level types,
predicates admitting proper classes, classes of set instances, and sets of
instances.
The $\Type$ is interpreted as $\TClass$.\\
Large types are used just for preventing primitive typing errors.\\
Medium types are for expressing the properties of proper classes.\\
Small types are expressing the proper classes, which contain as sets all interpretations of the terms of that type.\\
%(Another level may be here for expressing the proper classes, which contain only sets. It can be called "Small" while "Small" ones are renamed to "Tiny")\\
%Medium1 types are for expressing the properties of proper classes.
%Small(Maybe better call it "Tiny"?)
Tiny types are interpreted as sets, which contain as sets all interpretations of the terms of that type.\\

One class-valued candidate is:
\begin{definition}[dependent product, variant A]
\[
\Pi B = \cc{f\in \power \Sigma B }{ (f\functional) \land (\dom(f)=\Dom(B)) }
\]
The expression $\Dom(B)$ is used to include the case when $B$ is a proper family, not just a function.
\end{definition}
\begin{definition}[dependent product, variant B]
(Note: For sets the following definition is also okey:
\[
\Pi B = \cc{f\in \left(\union\ran(B)\right)^{\Dom(B)}}{\forall (x\in\dom(f)) ,\ (f(x) \in B(x))},
\]
but it is not working well for classes)

(This variant is intended for set-indexed, set-valued families; extending its size range requires a separate argument.)
\end{definition}

\begin{definition}[dependent product, variant C]
Another definition:\\
\[
\prod w \eqdef \cc{f}{\exists D,\ (f\family) \land (D = \dom(f)=\dom(w)) \land (\forall x\in D,\ f(x) \in w(x))}
\]
\end{definition}

\begin{definition}[sections of a family]
\[
\Gamma B\eqdef\cc{s:\Dom(B)\to\Sigma B}{(\pr_1\circ s) \subseteq \I}
\]
\end{definition}
%13
\begin{definition}[dependent product of a family]
\[
\Pi B \eqdef \cc{\pr_2\circ s}{s\in \Gamma B}
\]
\end{definition}

\begin{theorem}[Sections and dependent function graphs]
Let $A$ be a set and let $B(a)$ be a set for every $a\in A$.
Use injective set-pair coding, and put
\[
S=\{(a,b):a\in A,\ b\in B(a)\}.
\]
Then $S$ is a set, and there is a bijection between the set of functions
$f$ on $A$ with $f(a)\in B(a)$ and the set of sections
$s:A\to S$ of the projection $\pr_1:S\to A$. The mutually inverse maps are
\[
f\longmapsto(a\mapsto(a,f(a))),\qquad
s\longmapsto\pr_2\circ s.
\]
In particular the product is the set
\[
\{f\in\power S:f\text{ is functional and }\dom(f)=A\}.
\]
\end{theorem}
\Proof
Replacement collects the sets $\{a\}\times B(a)$, and union gives $S$.
A function graph contained in $S$ has exactly the required pointwise
values when its domain is $A$. Separation inside $\power S$ therefore
forms the product set. For such an $f$, replacement forms the graph of
$s(a)=(a,f(a))$, and $\pr_1\circ s=\operatorname{id}_A$.
Conversely, for a section $s$, write $s(a)=(c,b)\in S$.
The section equation gives $c=a$, hence $b\in B(a)$. Thus
$\pr_2\circ s$ is a member of the product. Applying the two constructions
in either order preserves every value, by the two projection equations
and injectivity of pairing. Function extensionality proves they are inverse.
\Done

This correspondence uses no choice axiom: it transforms supplied sections
and functions. It does not assert that the product is inhabited merely
because every $B(a)$ is nonempty. For arbitrary set-indexed families that
assertion is precisely the set axiom of choice. When $A=\emptyclass$,
both the product and the set of sections contain exactly the empty
function, so the correspondence includes the empty-index case.

\subsubsection{Soundness of $\Pi$--types}
\begin{theorem}[Soundness of $\Pi B$ Formation, Abstraction, Application, and $\beta$/$\eta$ Rules]
Let $B$ be a family of classes over a set domain $A=\Dom(B)$, and assume sethood of the graphs constructed below (as supplied by the class-function image principle). Let
\[
\Pi B = \cc{f \in \power \Sigma B}{(f\functional) \land (\dom(f) = \Dom(B)) \land (\forall (x \in \Dom(B)),\ (f(x) \in B(x)))}.
\]
Then:
\begin{enumerate}
\item \textbf{Introduction ($\lambda$-Abstraction):} If for every $(x \in \Dom(B))$ the term $t(x)$ is a set satisfying $(t(x) \in B(x))$, then the functional graph
\[
(\lambda (x \in \Dom(B)),\, t(x)) = \cc{(x, t(x))}{x \in \Dom(B)}
\]
is a functional subclass of $\Sigma B$ with domain $\Dom(B)$, hence $((\lambda (x \in \Dom(B)),\, t(x)) \in \Pi B)$.
\item \textbf{Elimination (Dependent Application):} If $(f \in \Pi B)$ and $(a \in \Dom(B))$, then the singleton fiber $\cc{y}{(a, y) \in f} = \{f(a)\}$ has unique element $f(a) = \mathrm{Value}(\cc{y}{(a, y) \in f}) = \bigcap\cc{y}{(a, y) \in f}$ satisfying $(f(a) \in B(a))$.
\item \textbf{$\beta$- and $\eta$-Computation:} For all $(a \in \Dom(B))$, $(((\lambda (x \in \Dom(B)),\, t(x)))(a) = t(a))$, and for every $(f \in \Pi B)$, $((\lambda (x \in \Dom(B)),\, f(x)) = f)$ by functional extensionality.
\end{enumerate}
These rules use the explicitly restricted product defined in this theorem. They do not follow from the unrestricted same-domain implication in the generated \texttt{DefPiPredicateModel} record discussed above.
\end{theorem}
\pagebreak
\subsection{Implementing W--types}
A W-type describes well-founded trees: a label $a$ determines the type
$B(a)$ of branches. The rules first state the intended interface. The
finite-stage model developed afterward is restricted to finite branching,
that is, finite $B(a)$ for each label $a$. This restriction is what allows
all branches of a node to fit in one finite stage. Infinite branching requires a further construction: after constructing
the M-type, we obtain the general set-arity W-type as its well-founded
part. It is not supplied by the finite-stage formulas alone.

\subsubsection{Axioms of W--types}
\underline{Formation of W--type:}\\
Let $B$ be a family of types.\\
Then $\W B$ is also a type.\\
\underline{Introduction of terms of W--type:}\\
\[
\sup(B):\prod_{(a\in \Dom(B))}((B(a)\to\mbox{W}B)\to \mbox{W}B)
\]
\underline{Elimination of terms of W--type:}\\
\[
\begin{aligned}
\mbox{W}_\mathrm{rect}:{}&
\prod_{(A : \Un)} \prod_{(B : A \to \Un)} \prod_{(P : \mathrm{W}B\to \Un)}\\
&\quad\left(\begin{aligned}
&\prod_{(a : A)} \prod_{(w : B(a) \to \mathrm{W}B)}\\
&\quad\left(\prod_{(f : B(a))} P(w(f))\right)\to P(\sup(B)(a)(w))
\end{aligned}\right)\\
&\quad\to\left(\prod_{(w : \mathrm{W}B)} P(w)\right)
\end{aligned}
\]
\underline{Computation of terms of W--type:}\\
\[
\mbox{W}_{\mathrm{rect}} (A,B,E,e,\sup(B,a,f))
\eqdef e(a, f, \lambda(b:B(a))\mapsto \mbox{W}_{\mathrm{rect}} (A,B,E,e, f(b)))
\]
\pagebreak
\subsubsection{Model of $\text{W}$--types}
% https://tex.stackexchange.com/questions/151549/mathscr-latex-or-xelatex
\underline{Formation of W--type:}\\

Assume $B$ is a family of classes with a class $A$ being its domain.
(a.k.a. ``some dependent class $B$, which depends on $\Dom(B)$'', where $\Dom(B)=A$.)

%\EuScript{W}
%\mathscr{W}
\[
(\EuScript{W}_{\mathrm{fun}}B)(0)\eqdef \emptyclass
\]
\[
(\EuScript{W}_{\mathrm{fun}}B)(n+1)\eqdef \cc{(r,s)}{ (r\in \Dom(B))\land (s:B(r)\to (\EuScript{W}_{\mathrm{fun}}B)(n) )}
\]
% why \textbf{N} instead of \mathbb{N}
\[
\EuScript{W}B \eqdef\bigcup_{n\in\omega}((\EuScript{W}_{\mathrm{fun}}B)(n)),
\]
Equivalently, when all stages are sets, this is the union of the range:
\[
\EuScript{W}B \eqdef\bigcup\ran(\EuScript{W}_{\mathrm{fun}}B).
\]
The range alone collects the stages, whereas its union collects their trees.

A variant of a definition, which works only for definable families and that uses %class--theoretic implementations of
the sigma types and the empty type:
\[
\mbox{W}_{\mathrm{fun}}':\text{Family.Expression}
\to
(
\mathbb{N}
%\text{.Expression}
\to
\Type
)
\]
\[
(\mbox{W}_{\mathrm{fun}}'B)(0)\equivdef \mathrm{Void}
\]
\[
(\mbox{W}_{\mathrm{fun}}'B)(n+1)\ \equivdef
\sum_{(r: \Dom(B))} (B(r)\to (\mbox{W}_{\mathrm{fun}}'B)(n) )
\]
%\Pi\Sigma\mathop{\dum}\limits_{(r\in \Dom(B))} \sum
%\mathop{\LARGE\ensuremath \mathrm{\LARGE \vphantom\sum W}\mathlarger{W}}\limits
\[
\mbox{W}':\text{Family.Expression}\to\Type
\]
\[
\mbox{W}'B\equivdef\bigsqcup_{n:\mathbb{N}} ((\mbox{W}_{\mathrm{fun}}'B)(n))
\]

For the finite-branching model, the intended fixed-point equation is
%=
\[
F(X)\cong X,
\]
where a functor $F$ is defined to be a polynomial functor $P_B$:
\[
P_B(X)
\cequiv
\sum_{a:A}
(B(a)
\to
X)
\]
\[
F(X)\cequiv P_B(X)
\]
For a polynomial functor $P$ with finite arities, define
\[
P_{\text{fun}}(0)\equivdef \text{Void}
\]
\[
P_{\text{fun}}(n+1)\equivdef P(P_{\text{fun}}(n))
\]
and the corresponding fixpoint type is
\[
\text{Fix}(P) \equivdef \bigsqcup_{n:\textbf{Nat}} P_{\text{fun}}(n).
\]
%$\big\Sigma$
%\scalebox{1.5}{$\Sigma$}
%\scalebox{2}{$\Sigma$}

%\vphantom{p}_B
\begin{lemma}[levels of W--types are cumulative]
For all $n\in\N$,
\[
(\mbox{W}_{\mathrm{fun}}B)(n)
\subseteq
(\mbox{W}_{\mathrm{fun}}B)(n+1).
\]
\Proof
By induction on $n$,\\
\underline{Base}: n=0.\\
 Trivially, $\emptyclass \subseteq (\mbox{W}_{\mathrm{fun}}B)(1)$.\\
\underline{Step of induction}:\\
Assume \[
(\mbox{W}_{\mathrm{fun}}B)(n)
\subseteq
(\mbox{W}_{\mathrm{fun}}B)(n+1).
\]
To show:
\[
(\mbox{W}_{\mathrm{fun}}B)(n+1)
\subseteq
(\mbox{W}_{\mathrm{fun}}B)(n+2).
\]
It follows from the assumption, that
\[
(B(r)\to(\mbox{W}_{\mathrm{fun}}B)(n))
\subseteq
(B(r)\to(\mbox{W}_{\mathrm{fun}}B)(n+1)).
\]
Therefore
\[
\sum_{(r\in \Dom(B))} (B(r)\to(\mbox{W}_{\mathrm{fun}}B)(n))
\subseteq
\sum_{(r\in \Dom(B))} (B(r)\to(\mbox{W}_{\mathrm{fun}}B)(n+1)).
\]
And finally,
\[
(\mbox{W}_{\mathrm{fun}}B)(n+1)
\subseteq
(\mbox{W}_{\mathrm{fun}}B)(n+2).
\]
\Done
\end{lemma}



\underline{Introduction of terms of W--type:}\\
\begin{definition}[constructor of W--type]
For fixed function $B$,
\[
\sup(B):\prod_{(a\in \Dom(B))}((B(a)\to\mbox{W}B)\to \mbox{W}B)
\]
Assume $(a\in \Dom(B))$ and $(s:B(a)\to\mbox{W}B)$. Since the branch set $B(a)$ is finite, its image $s[B(a)]$ is finite. If this image is nonempty, let $m$ be the maximum of the least stage indices of its elements; if it is empty, take $m=0$. Then $s[B(a)]\subseteq (\mbox{W}_{\mathrm{fun}}B)(m)$. Therefore $s:B(a)\to(\mbox{W}_{\mathrm{fun}}B)(m)$. Therefore $(a,s)\in (\mbox{W}_{\mathrm{fun}}B)(m+1)$, which is a subset of $\mbox{W}B$, so $(a,s)\in\mbox{W}B$. Therefore we can define $\sup$ this way:
\[
\sup(B) \eqdef ((a\in \Dom(B))\mapsto ((s:B(a)\to\mbox{W}B)\mapsto (a,s)) ).
\]
%Therefore $\sup(B) = \id(\mbox{W}B)$ and $\sup(B)\restriction  = \id(\mbox{W}B)$.
\end{definition}

\underline{Elimination of terms of W--type:}\\
\begin{definition}[destructor of W--type]
\[
\begin{aligned}
\mbox{W}_\mathrm{rect}:{}&
\prod_{(A : \Un)} \prod_{(B : A \to \Un)} \prod_{(P : \mathrm{W}B\to \Un)}\\
&\quad\left(\begin{aligned}
&\prod_{(a : A)} \prod_{(w : B(a) \to \mathrm{W}B)}\\
&\quad\left(\prod_{(f : B(a))} P(w(f))\right)\to P(\sup(B)(a)(w))
\end{aligned}\right)\\
&\quad\to\left(\prod_{(w : \mathrm{W}B)} P(w)\right)
\end{aligned}
\]
\begin{lstlisting}
W_rect
     : forall (A : Type) (B : A -> Type) (P : W B -> Type),
       (forall (a : A) (w : B a -> W B), (forall f : B a, P (w f)) -> P (sup B a w)) ->
       forall w : W B, P w
\end{lstlisting}
\end{definition}

\underline{Computation of terms of W--type:}\\

By finiteness of each branch set $B(a)$ and $(\mbox{W}_{\mathrm{fun}}B)(0)=\emptyclass$, it holds that
\[
\sup(B,a,f)\in (\mbox{W}_{\mathrm{fun}}B)(m+1)
\]
for some $m\in\N$.


\begin{theorem}[Dependent recursion for the finite-branching set model]
\label{thm:Wcomp}
Work in the usual set theory with powerset, replacement, and recursion on
$\omega$. Let $A$ be a set and each $B(a)$ a finite set. Write $W_n$ for
the stage above and $W=\bigcup_{n<\omega}W_n$, using ordinary injective
set-pair coding. Let $E(w)$ be a set for each $w\in W$. Suppose a
specified step operation assigns
\[
e(a,s,h)\in E((a,s))
\]
whenever $s:B(a)\to W$ and $h$ is a function with
$h(b)\in E(s(b))$ for every $b\in B(a)$. Then there is exactly one
function $r$ on $W$ with $r(w)\in E(w)$ and
\[
r((a,s))=e\bigl(a,s,b\mapsto r(s(b))\bigr).
\]
This is the interpretation of the displayed $\mathrm W_{\mathrm{rect}}$
and its computation rule in this set model.
\end{theorem}
\Proof
Powerset and replacement make every function set $W_n^{B(a)}$ and
hence each stage
\[
W_{n+1}=\sum_{a\in A}W_n^{B(a)}
\]
a set.
Recursion on $\omega$, replacement, and union give the set $W$.
The earlier cumulative-stage and finite-branching arguments show that
$W$ is closed under $(a,s)$ and that each node has a unique such
decomposition. In particular, a node in $W_{n+1}$ has every child in
$W_n$.

Put $r_0=\emptyclass$. Given $r_n$ on $W_n$, define on $W_{n+1}$
\[
r_{n+1}((a,s))=e\bigl(a,s,b\mapsto r_n(s(b))\bigr).
\]
The inductive typing assertion makes the right side defined and in
$E((a,s))$; replacement collects its graph. The crucial compatibility
assertion is $r_{n+1}|_{W_n}=r_n$. For $n=0$ it is vacuous. For the
step, a node $(a,s)\in W_{n+1}$ has children in $W_n$, where
$r_{n+1}$ and $r_n$ already agree. Thus its values under $r_{n+2}$
and $r_{n+1}$ are applications of the same $e$ to identical arguments.
This proves compatibility at the next stage.

The union $r=\bigcup_n r_n$ is therefore a function on $W$, with the
required typing. Every child family fits into some $W_n$, so the
stage equation gives the claimed computation equation. If $q$ is
another such function, induction on $n$ proves $q|_{W_n}=r_n$:
the base is empty and the step uses the computation equation on
children in $W_n$. Their unions agree, proving uniqueness.
\Done

\begin{corollary}[Structural induction and initiality]
If a predicate $Q$ on $W$ holds at $(a,s)$ whenever it holds at every
$s(b)$, then it holds everywhere on $W$. Furthermore, for every set $X$
and map $\alpha:P_B(X)\to X$ there is exactly one $r:W\to X$ such that
\[
r((a,s))=\alpha(a,r\circ s).
\]
Thus the constructor $P_B(W)\to W$ is an initial $P_B$-algebra.
\end{corollary}
\Proof
For induction, prove $Q$ on $W_n$ by induction on $n$; each node at the
next stage has all children at the preceding one. For initiality, use
the theorem with constant family $E(w)=X$ and
$e(a,s,h)=\alpha(a,h)$. Its uniqueness clause is exactly uniqueness of
the algebra morphism. The argument includes $X=\emptyclass$: a step
map into the empty set exists only when its domain is empty, and the
theorem then still applies.
\Done

\begin{remark}[Stages are identified along their inclusions]
In this cumulative set model, the notation $\bigsqcup_n W_n$ above
must mean the union, or equivalently the colimit along the inclusions
$W_n\subseteq W_{n+1}$. It cannot mean the disjoint union of the stages:
that would retain a separate tagged copy of a tree at each later stage,
and the proof would not identify their recursor values automatically.
The compatibility argument above supplies precisely the identification
needed by the model. A model with a proper class of labels or genuinely
class-valued dependent outputs requires additional representation and
recursion arguments and is not asserted by this set-model theorem.
\end{remark}

\begin{example}[Why the finite-stage proof needs finite branching]
Take distinct labels $z,u,c$, with $B(z)=\emptyclass$,
$B(u)=\{0\}$, and $B(c)=\omega$. The finite trees $t_0=(z,\emptyclass)$ and
$t_{n+1}=(u,\{(0,t_n)\})$ have unbounded finite heights.
Their sequence is a function $s:\omega\to W$, but its image is not
contained in any single $W_m$. Thus $(c,s)$ belongs to $P_B(W)$ but
not to the union of the finite stages. The missing root is well founded;
it merely has children of unbounded finite height. This shows failure
of this construction for infinite branching, not failure of the general
notion of a W-type.
\end{example}

\pagebreak
\subsection{Implementing M--types by finite paths}
An M-type allows infinite depth. The author's path-based proposal can be
made precise using finite addresses: a potentially infinite tree is
specified by all its finite observations. We now give a set-level model
and prove its universal property. No finite-branching hypothesis is needed.

Fix a set $A$ and a set family $(B(a))_{a\in A}$. Put
\[
C=\sum_{a\in A}B(a),\qquad
C^{<\omega}=\bigcup_{n<\omega} C^n.
\]
Elements of $C^{<\omega}$ are finite words; the empty word is $\epsilon$,
and concatenation is denoted by $pq$. A letter $(a,b)$ records the
parent label as well as its branch index, so overlapping branch sets do
not cause ambiguity.

\begin{definition}[Trees described by finite addresses]
Let $M_B$ consist of the partial functions $t:C^{<\omega}\rightharpoonup A$
with domain $D_t$ satisfying:
\begin{enumerate}
\item $\epsilon\in D_t$;
\item $D_t$ is prefix closed;
\item if $p\in D_t$ and $c\in C$, then
\[
pc\in D_t\iff c=(t(p),b)\text{ for some }b\in B(t(p)).
\]
\end{enumerate}
These are conditions on a subset of $C^{<\omega}\times A$, so powerset
and separation make $M_B$ a set. In particular $M_B=\emptyclass$ if
$A=\emptyclass$. Partial functions here are ordinary functional graphs;
their values are used only on their stated domains.
\end{definition}

For $t\in M_B$, let $a=t(\epsilon)$. For each $b\in B(a)$ define its
subtree $t_b$ by
\[
D_{t_b}=\{q:(a,b)q\in D_t\},\qquad t_b(q)=t((a,b)q).
\]
The tree axioms immediately give $t_b\in M_B$. Thus the destructor is
\[
\delta:M_B\longrightarrow P_B(M_B),\qquad
\delta(t)=(a,b\mapsto t_b).
\]

\begin{lemma}[Grafting and decomposition]
The destructor $\delta$ is bijective. Its inverse sends $(a,s)$, where
$s:B(a)\to M_B$, to the tree with root label $a$ and subtrees $s(b)$.
\end{lemma}
\Proof
Define $t$ on the set
\[
D_t=\{\epsilon\}\cup
\{(a,b)q:b\in B(a),\ q\in D_{s(b)}\}
\]
by $t(\epsilon)=a$ and $t((a,b)q)=s(b)(q)$.
Each nonempty address has a unique first letter and suffix, so this
is well defined. Prefix closure and the exact-child condition follow
at the root from the indexing set $B(a)$, and below it from the
corresponding conditions on each $s(b)$. Hence $t\in M_B$, with
$\delta(t)=(a,s)$. Conversely, decomposing any tree and grafting its
subtrees recovers every nonempty address and its label, as well as the
root. The constructions are inverse.
\Done

\begin{theorem}[Final coalgebra and corecursion]
\label{thm:M_final_coalgebra}
For every set $X$ and map $\gamma:X\to P_B(X)$ there is exactly one
$h:X\to M_B$ satisfying
\[
\delta\circ h=P_B(h)\circ\gamma,
\qquad P_B(h)(a,s)=(a,h\circ s).
\]
Thus $(M_B,\delta)$ is a final coalgebra for $P_B$ on sets.
Here $P_B$ is a functor: postcomposition preserves identities and
composition, so $P_B(\operatorname{id})=\operatorname{id}$ and
$P_B(k\circ h)=P_B(k)\circ P_B(h)$.
\end{theorem}
\Proof
Write $\gamma(x)=(a_x,s_x)$, where $s_x:B(a_x)\to X$.
Starting from a state $x$, define the state reached at a finite address
recursively on its length. The empty address reaches $x$. If $p$ reaches
$y$, then precisely the letters $(a_y,b)$ with $b\in B(a_y)$ are
admissible next, and $p(a_y,b)$ reaches $s_y(b)$. A word whose prefix
is inadmissible remains inadmissible. This is finite recursion with
unique transitions, and replacement collects the reached states on
each length; union over $\omega$ collects all valid addresses.

Label each valid address by $a_y$ for its reached state $y$. The result
$t_x$ satisfies the tree axioms by construction. Its root label is $a_x$,
and its subtree at $b$ is $t_{s_x(b)}$, because both use the identical
transition rule starting from $s_x(b)$. Thus $h(x)=t_x$ satisfies the
coalgebra equation. Replacement gives a set function $h$.

If $k:X\to M_B$ satisfies that equation, its root at $x$ has label
$a_x$ and its $b$-subtree is $k(s_x(b))$. Induction on address length
therefore proves that $k(x)$ and $t_x$ have exactly the same admissible
addresses of each length and the same labels there. Every address is
finite, so their graphs coincide. Consequently $k=h$.
The proof uses no selection of an infinite path and no choice function.
\Done

\begin{theorem}[Coinduction by bisimulation]
Let $R\subseteq M_B\times M_B$ have the following property: whenever
$t\mathrel R u$, their root labels agree, say with value $a$, and
$t_b\mathrel R u_b$ for every $b\in B(a)$. Then $t\mathrel R u$
implies $t=u$. In particular, equality is the largest such relation.
\end{theorem}
\Proof
For any related pair, the root labels agree. Inductively, equal labels
at a common address give exactly the same admissible next letters.
At each of those letters the corresponding subtrees are again related,
so their labels agree by the hypothesis. Induction on length proves
identical address domains and labels at every finite level. The trees
are their graphs on finite addresses, hence are equal. Equality itself
has the stated property, proving the final assertion.
\Done

\begin{example}[Streams and the difference from finite trees]
Let $A=S$ be a set of labels and $B(a)=\{0\}$ for every $a\in S$.
A tree then has exactly one address at each length, so reading successive
labels identifies $M_B$ with the stream set $S^\omega$. The destructor
reads the first label and the tail. If $S$ is nonempty, a constant stream
exists, while the finite-stage W-construction is empty: it has no
nullary label with which to start. If $S=\emptyclass$, both sets are
empty. Thus algebraic and coalgebraic fixed points have different
universal properties, even for the same polynomial functor.
\end{example}

\begin{proposition}[Finite trees inside the M-type]
Under the finite-branching assumptions of the W-model, there is a
canonical injective map $j:W\to M_B$. Its image consists exactly of
trees whose valid addresses have bounded finite length.
\end{proposition}
\Proof
Use the initiality of $W$ with the grafting algebra
$\delta^{-1}:P_B(M_B)\to M_B$. The resulting map satisfies
$j((a,s))=\delta^{-1}(a,j\circ s)$. Induction on the W-stage shows that
$j(w)$ has bounded address length.

To prove injectivity, use structural induction on $w$ with the stronger
assertion that for every $v\in W$, equality $j(w)=j(v)$ implies $w=v$.
Write $w=(a,s)$ and $v=(d,t)$. Applying $\delta$ to the equality gives
$a=d$ and $j(s(b))=j(t(b))$ for each $b\in B(a)$. The induction
hypotheses for the children $s(b)$ apply to these arbitrary children
$t(b)$, giving $s(b)=t(b)$. Function extensionality then gives $s=t$
and hence $w=v$.

Conversely, induct on an upper bound $n$ for the address lengths of a
tree $u$. If $n=0$, the root has no children, so its label $a$ has
$B(a)=\emptyclass$ and $u=j((a,\emptyclass))$. If $n>0$, every root
subtree has address lengths at most $n-1$, hence is $j(w_b)$ for some
$w_b\in W$. Injectivity makes each $w_b$ unique, so replacement gives
the function $b\mapsto w_b$ without a choice axiom. Constructor closure
of $W$ gives the node $(a,b\mapsto w_b)$, and its image under $j$ is
$u$ by grafting. This completes the induction.
\Done

\subsection{W-types with arbitrary set arities}
The finite-stage construction fails for some infinite arities, but the
M-type supplies a set in which to isolate the well-founded trees. Keep
$A$ and every $B(a)$ sets; no finiteness assumption is made here.

\begin{definition}[The well-founded part of the M-type]
For $t\in M_B$, put $q\prec_t p$ when $q$ is an immediate child address
of $p$ in $D_t$. Say that $t$ is well founded if every nonempty subset
of $D_t$ has a $\prec_t$-minimal member. Define
\[
W_B^{\mathrm{wf}}=\{t\in M_B:\prec_t\text{ is well founded}\}.
\]
Separation makes this a set. Well-foundedness here is the minimal-member
condition, not merely an assertion that no infinite path has been chosen.
\end{definition}

\begin{lemma}[Grafting preserves and detects well-foundedness]
For $t=\delta^{-1}(a,s)$, one has
\[
t\in W_B^{\mathrm{wf}}
\iff s(b)\in W_B^{\mathrm{wf}}\text{ for every }b\in B(a).
\]
Consequently grafting restricts to a bijection
$P_B(W_B^{\mathrm{wf}})\to W_B^{\mathrm{wf}}$.
\end{lemma}
\Proof
A nonempty subset of addresses of $s(b)$ lifts by prefixing $(a,b)$
to a nonempty subset of $D_t$. A minimal member of the latter gives a
minimal member of the former, since immediate children stay in that
subtree. This proves the forward direction.

Conversely, let $S\subseteq D_t$ be nonempty. If it contains only the
root, that root is minimal. Otherwise choose one nonroot address in
$S$, and let $(a,b)$ be its first letter. The suffixes of the addresses
of $S$ beginning with this letter form a nonempty subset of $D_{s(b)}$.
Choose a minimal suffix there. Its prefixed address is minimal in all
of $S$, because every immediate child has the same first letter.
Thus $t$ is well founded. The bijection now follows from the already
proved grafting/decomposition bijection on $M_B$.
\Done

\begin{lemma}[Ordinal heights and the subtree relation]
Every $t\in W_B^{\mathrm{wf}}$ has a uniquely determined ordinal height
$h(t)$, with
\[
h(t)=\sup_{b\in B(t(\epsilon))}\bigl(h(t_b)+1\bigr).
\]
The relation $u\prec t$ on $W_B^{\mathrm{wf}}$ meaning that $u$ is an
immediate root subtree of $t$ is set-like and well founded.
\end{lemma}
\Proof
Apply the ordinal-rank construction for a founded set relation to
$(D_t,\prec_t)$. It assigns ordinals $r_t(p)$ satisfying
\[
r_t(p)=\sup_{q\prec_t p}(r_t(q)+1).
\]
All predecessor collections are sets, so replacement and ordinal union
supply the suprema. Define $h(t)=r_t(\epsilon)$. The suffix bijection
between $D_{t_b}$ and the addresses of $t$ beginning with
$(t(\epsilon),b)$ preserves the child relation. Restriction of $r_t$
therefore satisfies exactly the rank recursion for $t_b$; uniqueness
of founded recursion gives
$h(t_b)=r_t((t(\epsilon),b))$. This proves the displayed height formula,
and in particular $h(t_b)<h(t)$.

The immediate subtrees form the image of the set $B(t(\epsilon))$,
so the subtree relation is set-like. For any nonempty subset $S$ of
$W_B^{\mathrm{wf}}$, replacement collects the set of its heights.
Choose its least ordinal and a tree in $S$ of that height. A predecessor
in $S$ would have smaller height, which is impossible. Hence the relation
is well founded. This argument constructs heights from founded address
relations first; it does not assume foundedness of the subtree relation.
Finally, if another ordinal assignment satisfies the displayed height
equation, founded induction on this now-established subtree relation
makes it agree with $h$ at every tree: agreement on all immediate
subtrees makes the two suprema identical.
\Done

\begin{theorem}[Dependent recursion for arbitrary set arities]
\label{thm:general_W_recursion}
Let $(E(t))_{t\in W_B^{\mathrm{wf}}}$ be a specified family of sets. Suppose a
specified operation assigns
\[
e(a,s,k)\in E(\delta^{-1}(a,s))
\]
whenever $s:B(a)\to W_B^{\mathrm{wf}}$ and $k$ is a function on $B(a)$
with $k(b)\in E(s(b))$ for every $b\in B(a)$. There is exactly one function
$r$ with domain $W_B^{\mathrm{wf}}$, values $r(t)\in E(t)$, and
\[
r(\delta^{-1}(a,s))=e(a,s,b\mapsto r(s(b))).
\]
\end{theorem}
\Proof
Use Theorem~\ref{thm:setlike_class_recursion} on the founded immediate-
subtree relation. Given a function $g$ on the predecessors of
$t=\delta^{-1}(a,s)$, form $k(b)=g(s(b))$. If every such value belongs
to $E(s(b))$, use $e(a,s,k)$ as the recursion step; otherwise use the
empty set as a temporary default. For any other set input, including
a nonfunction or a function whose domain is not the predecessor set,
also return the empty set. This defines a total set-valued step,
so founded recursion gives a function $r$.
Founded induction shows that its values are correctly typed: assuming
$r(u)\in E(u)$ at every predecessor, the first branch of the step applies
and yields $r(t)\in E(t)$. Thus the temporary default never occurs,
and the required equation holds everywhere. Replacement makes the graph
a set. Two typed solutions agree on each tree as soon as they agree
on its predecessors; founded induction therefore proves uniqueness.
Repeated identical subtrees cause no ambiguity, since $k$ uses the same
value of $g$ at each occurrence.
\Done

\begin{corollary}[Structural induction and the general initial algebra]
A predicate on $W_B^{\mathrm{wf}}$ that is preserved by grafting from
all immediate subtrees holds everywhere. The grafting map makes
$W_B^{\mathrm{wf}}$ an initial algebra for $P_B$ on sets.
\end{corollary}
\Proof
The induction assertion is founded induction for the subtree relation.
Given any algebra $\alpha:P_B(X)\to X$, apply the recursion theorem
with constant family $E(t)=X$ and $e(a,s,k)=\alpha(a,k)$.
The resulting unique $r$ satisfies
$r\circ\delta^{-1}=\alpha\circ P_B(r)$, exactly the algebra-morphism
condition. This includes the empty cases: existence of the specified
step operation is part of the hypothesis.
\Done

\begin{proposition}[Agreement with the finite-stage model]
If every $B(a)$ is finite, the canonical map $j:W\to M_B$ constructed
above identifies $W$ with $W_B^{\mathrm{wf}}$.
\end{proposition}
\Proof
A tree with address lengths bounded by a natural number is founded:
in any nonempty subset of its addresses choose one of greatest length;
none of its immediate children can belong to that subset. Conversely,
apply founded induction to a tree in $W_B^{\mathrm{wf}}$ and the
property of having bounded finite address lengths. If every immediate
subtree has that property, let $n_b$ be its least natural-number bound.
These bounds are uniquely specified, so replacement collects them.
There are finitely many children; taking the maximum of their bounds
and adding one bounds the whole tree. At a leaf the bound is zero.
Thus every founded tree has bounded finite height. The preceding
characterization of the image of $j$ now proves the assertion.
No equivalence between pathlessness and foundedness, and no selection of
an infinite branch, is used.
\Done

The earlier infinitely branching root with children of all finite heights
belongs to $W_B^{\mathrm{wf}}$ although it is absent from the finite-stage
union. Its address relation is founded by the grafting lemma, and its
ordinal height is $\omega$. This is precisely the extra behavior supplied
by arbitrary set arities.

\paragraph{Relation to the larger class-valued proposal.}
The original sketch described collections of paths with finite or
countably infinite domains and a compatibility condition on common
prefixes. The finite-address construction supplies that compatibility
and the destructor/corecursion rules for set labels and set arities.
It does not construct a set of all trees when the labels or branch
collections are proper classes. That larger proposal still requires an
explicit ambient size convention and a representation theorem; its
universal property cannot be inferred merely from the set case.
\ifdefined\qwertyflag
\undef\qwertyflag
\end{document}
\fi
